\chapter{Algebraic structure of mixed-state renormalization fixed points}%
\label{ch:mpdo_rfp_algebraic}

In the general mixed-state renormalization analysis of
\cite[Section~4]{Cirac2016MPDO_arXiv}, an MPDO tensor in vertical canonical form
generates closed-chain operators whose algebraic structure is governed by
diagonal $\chi$-matrices and fusion isometries $U_{\alpha,\beta}$. This
chapter treats the algebras of blocked fixed points, the structure
coefficients $\tr(\chi^L)$ and the idempotent condition. It formulates the
algebra conditions for a basis of normal tensors, shows that under these
conditions the blocked coefficients are traces of powers of positive diagonal
matrices
(Corollaries~\ref{cor:mpdo_bnt_label_theorem_data_blocked_coeff_eq_positive_blocked_chi_trace_pow}
and~\ref{cor:mpdo_bnt_label_theorem_witness_blocked_coeff_eq_positive_blocked_chi_trace_pow}),
and normalizes the invertible gauges between two-site and one-site sectors
to unitary conjugations. The resulting physical channels prove that the
algebra conditions make the tensor a renormalization fixed point
(Lemma~\ref{thm:mpdo_bnt_algebra_tensor_to_rfp}), one implication of
\cite[Theorem~4.14]{Cirac2016MPDO_arXiv}. The next chapter gives fixed points
whose structure coefficients depend on the length, and the chapter after it
treats the fusion isometries.

\section{Closed-chain operators}

As established in Lemma~\ref{thm:vertical_cf_of_horizontal_cf} following
\cite[Section~4]{Cirac2016MPDO_arXiv}, the vertical canonical form of an MPDO
tensor is $\widetilde M = U^\dagger B U$ with
$B=\bigoplus_\alpha\mu_\alpha\otimes M_\alpha$.

The closed-chain operators $O_L(M_\alpha)$ contract cyclically the BNT bond
indices of $L$ copies in the vertical reading. Decomposing each vertical
physical index into a pair $(a,b)$ of horizontal bond indices and writing
$\widehat M_\alpha^{ab}=M_\alpha^{(a,b)}$, the matrix entries for index strings
$\mathbf a=(a_1,\ldots,a_L)$ and $\mathbf b=(b_1,\ldots,b_L)$ are
\begin{align}
    [O_L(M_\alpha)]_{\mathbf a,\mathbf b}
     & = \tr\!\left(\widehat M_\alpha^{a_1b_1}\cdots
    \widehat M_\alpha^{a_Lb_L}\right) \notag                    \\
     & = [\rho^{(L)}(\widehat M_\alpha)]_{\mathbf a,\mathbf b},
    \notag
\end{align}
where the trace runs over the BNT bond space. In the operator algebra of
\cite[Theorem~4.14]{Cirac2016MPDO_arXiv}, these closed-chain operators
$O_L(M_\alpha)=\rho^{(L)}(\widehat M_\alpha)$ act on the horizontal bond
strings $\mathbf a$ and $\mathbf b$.

\section{Algebra structure}

% Source-numbering note: arXiv:1606.00608 uses one theorem counter for
% definitions, propositions, theorems, examples, and observations in Section 4.
% The three examples after Theorem 4.9 are numbered 4.10--4.12; hence
% `thm:IV.13` is Theorem 4.14 in the arXiv source.

\begin{definition}[Support algebras and blocking maps]%
    \label{def:mpdo_algebra_structure_data}
    \lean{MPOTensor.AlgebraStructureData}
    \leanok
    \uses{def:mpo_tensor}
    Let $\{\mathcal{A}_n\}_{n \ge 0}$ be a family of $*$-subalgebras of
    $\MN{D}$, one at every blocking size, equipped with bilinear maps
    $m_n : \mathcal{A}_n \times \mathcal{A}_n \to \mathcal{A}_{2n}$ and
    linear inclusion maps $\iota_n : \mathcal{A}_n \to \mathcal{A}_{n+1}$.
    These maps are realized in the ambient bond-space matrix algebra by
    \begin{align}
        m_n(x,y)   & = xy, \notag \\
        \iota_n(x) & = x.
        \notag
    \end{align}
\end{definition}

\begin{definition}[Algebras of blocked adjoint fixed points]%
    \label{def:mpdo_blocked_fixed_point_algebra_tower}
    \lean{MPOTensor.HasBlockedAdjointFixedPointAlgebraTower}
    \leanok
    \uses{def:mpdo_algebra_structure_data, def:mpo_tensor}
    An MPO tensor $M$ has a family of algebras of blocked adjoint fixed points
    if there exists a family of support algebras with blocking and inclusion
    maps as in Definition~\ref{def:mpdo_algebra_structure_data} such that, for
    every positive blocking size $n$,
    \begin{align}
        \mathcal{A}_n & = \Fix(E_n^{\dagger}),
        \notag
    \end{align}
    where $E_n$ is the blocked transfer map. This equality is termed
    compatibility with $M$. This condition concerns adjoint fixed points;
    it is not the algebra condition of
    \cite[Theorem~4.14(ii)]{Cirac2016MPDO_arXiv}, which requires a common
    BNT $\{M_\alpha\}_\alpha$, the structure coefficients
    $c_{\alpha,\beta,\gamma}^{(L)}=\tr(\chi_{\alpha,\beta,\gamma}^{L})$,
    and a vector $m$, with $m_\alpha=\tr(\mu_\alpha)$, that is an idempotent
    for the multiplication induced by $c^{(1)}$.
\end{definition}

\begin{theorem}[Idempotent transfer maps and stationary support algebras]%
    \label{thm:mpdo_blocked_fixed_point_algebra_tower_from_zcl}
    \lean{MPOTensor.hasBlockedAdjointFixedPointAlgebraTower_of_isZCL_of_isTP_of_posDef_fixed}
    \leanok
    \uses{def:mpo_is_zcl, def:mpdo_blocked_fixed_point_algebra_tower}
    Assume the transfer map of the doubled tensor is trace-preserving and
    admits a positive-definite fixed point. If $E_1^2 = E_1$, then there is a
    stationary family of support algebras
    \begin{align}
        \mathcal{A}_n & = \Fix(E_1^{\dagger}),
        \notag
    \end{align}
    with multiplication $m_n(x,y) = xy$ and inclusion $\iota_n(x) = x$.
\end{theorem}
\begin{proof}\leanok
    Under the trace-preserving normalization and the faithful fixed point, Wolf's
    fixed-point-algebra theorem gives the $*$-subalgebra
    $\Fix(E_1^{\dagger}) \subseteq \MN{D}$. Since $E_1$ is
    idempotent, $E_n = E_1^n = E_1$ for every $n \ge 1$, so
    $\Fix(E_n^{\dagger}) = \Fix(E_1^{\dagger})$ at every
    positive blocking size, giving the stationary family
    $\mathcal{A}_n = \Fix(E_1^{\dagger})$.
\end{proof}

For the one-site and two-site vertical canonical forms, let
$\{M_\alpha\}_\alpha$ and $\{A_\alpha\}_\alpha$ denote their respective
BNTs. Their boundary contractions have the same local form:
\begin{center}
    \tenkzkernel
    % Formula: M_\alpha(X)_{ij}
    % = \sum_{a,b}(M_\alpha^{ij})_{ab}X_{ba}
    % = \tr(M_\alpha^{ij}X).
    % Source verdict: Papers/1606.00608/DAVID_Fig10.png, cited at source TeX
    % lines 1974--1980, places X immediately left of the boxed M_\alpha.
    % Ink-to-index and contraction audit: the direct adjacent virtual bond and
    % the single-row periodic return below together contract a and b;
    % the unlabeled upper and lower physical stubs carry the free i and j.
    % Paired signature (formula and network):
    % boundary|virtual-west=0|virtual-east=0|physical-up=1|physical-down=1.
    $M_\alpha(X)=$
    \begin{tenkz}[rows={wire}, cols=2, boundary=periodic]
        \tn[skin=box]{X} &
        \tn[skin=mpo, label pos=45, ports={90:physical, 270:physical}]{M_\alpha}
    \end{tenkz}
    \qquad
    % Formula: A_\alpha(X)_{ij}
    % = \sum_{a,b}(A_\alpha^{ij})_{ab}X_{ba}
    % = \tr(A_\alpha^{ij}X).
    % Source verdict: the second panel of DAVID_Fig10.png has the same geometry,
    % with X immediately left of the boxed A_\alpha.
    % Ink-to-index and contraction audit: the direct adjacent virtual bond and
    % the single-row periodic return below together contract a and b;
    % the unlabeled upper and lower physical stubs carry the free i and j.
    % Paired signature (formula and network):
    % boundary|virtual-west=0|virtual-east=0|physical-up=1|physical-down=1.
    $A_\alpha(X)=$
    \begin{tenkz}[rows={wire}, cols=2, boundary=periodic]
        \tn[skin=box]{X} &
        \tn[skin=mpo, label pos=45, ports={90:physical, 270:physical}]{A_\alpha}
    \end{tenkz}
\end{center}
In each cell, $X$ is inserted on the virtual bond and the virtual indices are
closed by $\tr$; the physical legs remain open. In
\cite[Appendix~C.4, lines~1974--1980]{Cirac2016MPDO_arXiv},
the weighted direct sums of these contractions are fixed by compositions
of the reduced tpCPM. Products of these fixed operators are then used to
prove that the compositions are identities on the corresponding direct sums
of matrix algebras. This argument concerns the vertical canonical forms,
not the fixed-point algebras $\Fix(E_n^\dagger)$ of
Definition~\ref{def:mpdo_blocked_fixed_point_algebra_tower}.

\begin{theorem}[Stabilized blocked adjoint fixed points yield support algebras]%
    \label{thm:mpdo_algebra_from_stabilized_adjoint_fixedpoints}
    \lean{MPOTensor.hasBlockedAdjointFixedPointAlgebraTower_of_adjointFixedPoints_eq_of_isTP_of_posDef_fixed}
    \leanok
    \uses{def:mpdo_blocked_fixed_point_algebra_tower}
    Assume the transfer map of the doubled tensor is trace-preserving and
    admits a positive-definite fixed point. If, for every positive blocking
    size $n$,
    \begin{align}
        \Fix(E_n^{\dagger}) & = \Fix(E_1^{\dagger}),
        \notag
    \end{align}
    then there is a compatible family of support algebras.
\end{theorem}
\begin{proof}\leanok
    The trace-preserving map and its positive-definite fixed point give a
    $*$-subalgebra $\Fix(E_1^\dagger)$, and setting
    $\mathcal{A}_n := \Fix(E_1^{\dagger})$ for every $n$ with ambient
    multiplication and inclusion yields
    $\mathcal{A}_n = \Fix(E_n^{\dagger})$ at every positive blocking
    size $n$ by the stabilization hypothesis
    $\Fix(E_n^{\dagger}) = \Fix(E_1^{\dagger})$, which is compatibility with $M$.
\end{proof}

Stabilization of adjoint fixed points does not establish the converse direction of
\cite[Theorem~4.14]{Cirac2016MPDO_arXiv}. In particular,
\cite[Theorem~4.14(ii)]{Cirac2016MPDO_arXiv} requires positive diagonal
matrices $\chi_{\alpha,\beta,\gamma}$ for the multiplication of the vertical
operators $O_L(M_\alpha)$ and a vector $m_\alpha=\tr(\mu_\alpha)$ that is
an idempotent for the multiplication induced by $c^{(1)}$.
By contrast, choosing a basis in each $\mathcal{A}_n$ gives coordinates for
multiplication and inclusion between different blocking sizes. These
coordinates alone do not identify the vertical BNT or its structure
coefficients.

\begin{definition}[Coordinates in the support algebras]%
    \label{def:mpdo_blocked_coefficients}
    \lean{MPOTensor.AlgebraStructureData.toBlockedCoefficients}
    \leanok
    \uses{def:mpdo_algebra_structure_data}
    For each blocking size $n$, fix a basis
    $\{b_i^{(n)}\}_{i \in I_n}$ of $\mathcal{A}_n$ over a finite index set
    $I_n$. The coordinates of $x \in \mathcal{A}_n$ are the unique scalars
    $(x_i)_{i \in I_n} \in \C^{I_n}$ such that
    \begin{align}
        x & = \sum_{i \in I_n} x_i\, b_i^{(n)}.
        \notag
    \end{align}
\end{definition}

\begin{theorem}[Reconstruction from coordinates]%
    \label{thm:mpdo_reconstruct_blocked_coefficients}
    \lean{MPOTensor.AlgebraStructureData.reconstructFromBlockedCoefficients_apply}
    \leanok
    \uses{def:mpdo_blocked_coefficients}
    Reconstructing a coordinate family $(a_i)_{i\in I_n}$ gives
    $\sum_i a_i b_i^{(n)}$ in $\mathcal{A}_n$.
\end{theorem}
\begin{proof}\leanok
    This is the finite-basis reconstruction formula for the chosen basis of
    $\mathcal{A}_n$.
\end{proof}

\begin{definition}[Multiplication coefficients between blocking sizes]%
    \label{def:mpdo_blocked_structure_coefficients}
    \lean{MPOTensor.AlgebraStructureData.blockedStructureCoefficients}
    \leanok
    \uses{def:mpdo_blocked_coefficients}
    For $i,j \in I_n$, define the multiplication coefficients
    $(c^{(n)}_{i,j,k})_{k \in I_{2n}}$ by
    \begin{align}
        m_n(b_i^{(n)},b_j^{(n)})
         & = \sum_{k \in I_{2n}} c^{(n)}_{i,j,k}\, b_k^{(2n)}.
        \notag
    \end{align}
    These are coordinates in $\mathcal{A}_{2n}$, not the same-length
    structure coefficients of the vertical operators $O_L(M_\alpha)$.
\end{definition}

\begin{theorem}[Reconstruction of multiplication from its coefficients]%
    \label{thm:mpdo_mul_coefficients_reconstruct}
    \lean{MPOTensor.AlgebraStructureData.coe_mul_eq_sum_blockedStructureCoefficients}
    \leanok
    \uses{def:mpdo_blocked_structure_coefficients}
    In the ambient bond-space matrix algebra,
    $b_i^{(n)} b_j^{(n)}
        = \sum_{k \in I_{2n}} c^{(n)}_{i,j,k} b_k^{(2n)}$.
\end{theorem}
\begin{proof}\leanok
    Applying the reconstruction formula of
    Theorem~\ref{thm:mpdo_reconstruct_blocked_coefficients} to the element
    $m_n(b_i^{(n)},b_j^{(n)}) \in \mathcal{A}_{2n}$ yields the expansion.
\end{proof}

\begin{definition}[Inclusion coefficients between blocking sizes]%
    \label{def:mpdo_blocked_inclusion_coefficients}
    \lean{MPOTensor.AlgebraStructureData.blockedInclusionCoefficients}
    \leanok
    \uses{def:mpdo_blocked_coefficients}
    For $i \in I_n$, define the inclusion coefficients
    $(d^{(n)}_{ij})_{j \in I_{n+1}}$ by
    \begin{align}
        \iota_n(b_i^{(n)})
         & = \sum_{j \in I_{n+1}} d^{(n)}_{ij}\, b_j^{(n+1)}.
        \notag
    \end{align}
\end{definition}

\begin{theorem}[Reconstruction of inclusion from its coefficients]%
    \label{thm:mpdo_inclusion_coefficients_reconstruct}
    \lean{MPOTensor.AlgebraStructureData.reconstructFromBlockedInclusionCoefficients}
    \leanok
    \uses{def:mpdo_blocked_inclusion_coefficients}
    Reconstructing the inclusion coefficients recovers the image of
    $b_i^{(n)}$ under $\iota_n$:
    $\sum_{j \in I_{n+1}} d^{(n)}_{ij} b_j^{(n+1)}
        = \iota_n(b_i^{(n)})$.
\end{theorem}
\begin{proof}\leanok
    \uses{def:mpdo_blocked_coefficients}
    Coefficient extraction and reconstruction are mutually inverse coordinate
    maps for the chosen basis of $\mathcal{A}_{n+1}$. Applying them successively
    to $\iota_n(b_i^{(n)})$ therefore returns the same element.
\end{proof}

\begin{theorem}[Reconstruction of compatible adjoint fixed points]%
    \label{thm:mpdo_fixedpoint_coefficients_reconstruct}
    \lean{MPOTensor.AlgebraStructureData.reconstructFromBlockedCoefficients_of_fixedPoint}
    \leanok
    \uses{def:mpdo_blocked_fixed_point_algebra_tower, def:mpdo_blocked_coefficients}
    If the support algebras are compatible with an MPO tensor $M$ and
    $E_n^\dagger X = X$ for a positive blocking size $n$, then
    $X \in \mathcal{A}_n$. Writing $(X_i)_{i \in I_n}$ for its coordinates
    gives $X = \sum_{i \in I_n} X_i b_i^{(n)}$.
\end{theorem}
\begin{proof}\leanok
    Compatibility identifies $\mathcal{A}_n$ with the adjoint fixed-point algebra of
    $E_n$, so $X$ belongs to $\mathcal{A}_n$.
    Theorem~\ref{thm:mpdo_reconstruct_blocked_coefficients} then reconstructs it.
\end{proof}

\begin{theorem}[Descent of blocked adjoint fixed points]%
    \label{thm:mpdo_adjoint_fix_descent}
    \lean{MPOTensor.adjoint_transferMap_apply_of_hasBlockedAdjointFixedPointAlgebraTower}
    \leanok
    \uses{def:mpdo_blocked_fixed_point_algebra_tower}
    If an MPO tensor has a compatible family of support algebras, then for
    every positive blocking size $n$ and every $X$,
    $E_n^{\dagger}X=X$ implies $E_1^{\dagger}X=X$.
\end{theorem}
\begin{proof}\leanok
    Compatibility at size $n$ places $X$ in $\mathcal{A}_n$. The inclusion
    $\iota_n$ places the same ambient matrix $X$ in $\mathcal{A}_{n+1}$.
    Compatibility at size $n+1$ therefore gives $E_{n+1}^\dagger X=X$.
    Since $E_{n+1}=E_n\circ E_1$, it follows that
    $X=E_{n+1}^\dagger X=E_1^\dagger E_n^\dagger X=E_1^\dagger X$.
\end{proof}

\begin{lemma}[Reverse inclusion of adjoint fixed points]%
    \label{lem:mpdo_adjoint_fix_reverse}
    \lean{MPOTensor.adjoint_blockedTransferMap_apply_of_adjoint_transferMap_apply}
    \leanok
    For every $n$, $E_1^{\dagger}X=X$ implies $E_n^{\dagger}X=X$.
\end{lemma}
\begin{proof}\leanok
    Without requiring a compatible family of support algebras, induction on
    $n$ with base case $E_0 = \id$, using $E_{n+1} = E_n \circ E_1$ and
    $(A \circ B)^{\dagger} = B^{\dagger} \circ A^{\dagger}$, proves the
    claim.
\end{proof}

\begin{lemma}[Blocked powers of adjoint eigenvectors]%
    \label{lem:mpdo_adjoint_blocked_eigenvector}
    \lean{MPOTensor.adjoint_blockedTransferMap_apply_of_adjoint_transferMap_eigenvector}
    \leanok
    If $E=E_1$ and $E^\dagger X=\lambda X$, then
    $E_n^\dagger X=\lambda^n X$ for every $n$.
\end{lemma}
\begin{proof}\leanok
    The case $n=0$ is $E_0=\id$. For the induction step,
    $E_{n+1}=E_n\circ E_1$, hence
    \begin{align}
        E_{n+1}^{\dagger}X
         & = E_1^{\dagger}E_n^{\dagger}X
        = E_1^{\dagger}(\lambda^n X)
        = \lambda^n E_1^{\dagger}X
        = \lambda^{n+1}X.
        \notag
    \end{align}
\end{proof}

\begin{theorem}[Finite-order adjoint eigenvalues under compatibility]%
    \label{thm:mpdo_algebra_root_of_unity_eigenvalue}
    \lean{MPOTensor.adjoint_transferMap_eigenvalue_eq_one_of_hasBlockedAdjointFixedPointAlgebraTower}
    \leanok
    \uses{def:mpdo_blocked_fixed_point_algebra_tower, thm:mpdo_adjoint_fix_descent, lem:mpdo_adjoint_blocked_eigenvector}
    Assume there is a compatible family of support algebras. If $X\ne 0$,
    $E_1^\dagger X=\lambda X$, and $\lambda^n=1$ for some $n>0$, then
    $\lambda=1$.
\end{theorem}
\begin{proof}\leanok
    By Lemma~\ref{lem:mpdo_adjoint_blocked_eigenvector},
    $E_n^\dagger X=\lambda^n X=X$. Theorem~\ref{thm:mpdo_adjoint_fix_descent} gives
    $E_1^\dagger X=X$. Thus $\lambda X=X$, so $(\lambda-1)X=0$. Since
    $X\ne0$, one obtains $\lambda=1$.
\end{proof}

\begin{theorem}[Equality of compatible adjoint fixed-point algebras]%
    \label{thm:mpdo_adjoint_fix_equality_from_algebra}
    \lean{MPOTensor.adjoint_blockedTransferMap_apply_iff_of_hasBlockedAdjointFixedPointAlgebraTower}
    \leanok
    \uses{def:mpdo_blocked_fixed_point_algebra_tower, thm:mpdo_adjoint_fix_descent, lem:mpdo_adjoint_fix_reverse}
    If an MPO tensor has a compatible family of support algebras, then
    $\Fix(E_n^{\dagger})=\Fix(E_1^{\dagger})$ for every positive blocking
    size $n$.
\end{theorem}
\begin{proof}\leanok
    Combining Theorem~\ref{thm:mpdo_adjoint_fix_descent} with
    Lemma~\ref{lem:mpdo_adjoint_fix_reverse} gives both inclusions.
\end{proof}

\begin{theorem}[Compatibility of the stationary family of support algebras]%
    \label{thm:mpdo_stationary_tower_from_algebra}
    \lean{MPOTensor.stationaryOfFaithfulFixedPoint_compatible_of_hasBlockedAdjointFixedPointAlgebraTower}
    \leanok
    \uses{def:mpdo_blocked_fixed_point_algebra_tower, thm:mpdo_adjoint_fix_equality_from_algebra}
    Under the trace-preserving normalization and a positive-definite fixed
    point, if $M$ has a compatible family of support algebras, then the
    stationary family $\mathcal{A}_n = \Fix(E_1^\dagger)$ is
    itself compatible with $M$: for every positive blocking size $n$,
    $\Fix(E_1^\dagger)=\Fix(E_n^\dagger)$.
\end{theorem}

\begin{theorem}[Compatible support algebras do not imply an idempotent transfer map]%
    \label{thm:mpdo_algebra_tower_not_imply_zcl}
    \lean{MPOTensor.exists_hasBlockedAdjointFixedPointAlgebraTower_not_isZCL}
    \leanok
    \uses{def:mpdo_blocked_fixed_point_algebra_tower, def:mpo_is_zcl, thm:mpdo_algebra_from_stabilized_adjoint_fixedpoints}
    There exists an MPO tensor of physical and bond dimension two whose
    doubled-tensor transfer map is trace-preserving and has a
    positive-definite fixed point, and which has a compatible family of
    support algebras, although its doubled-tensor transfer map is not
    idempotent.
\end{theorem}
\begin{proof}\leanok
    For the phase-flip channel with Kraus operators
    $K_0=\frac35 \Id$ and $K_1=\frac45 Z$, where $Z=\diag(1,-1)$,
    the transfer map satisfies
    \begin{align}
        E(X)_{ij}
         & =
        \begin{cases}
            X_{ij},            & i=j,     \\
            -\frac7{25}X_{ij}, & i\neq j.
        \end{cases}
        \notag
    \end{align}
    The identity is a positive-definite fixed point, and
    $K_0^\dagger K_0+K_1^\dagger K_1=\Id$ proves trace preservation.
    For every $n>0$ and $i\neq j$,
    $E^n(X)_{ij}=(-7/25)^n X_{ij}$. Since $(-7/25)^n\neq 1$ for every $n>0$
    and $E$ is self-adjoint,
    $\Fix((E^n)^\dagger)=\Fix(E^\dagger)$. Thus a compatible family of support
    algebras exists. On the other hand, for $X_{01}\ne0$,
    $E^2(X)_{01}=\frac{49}{625}X_{01}\neq-\frac7{25}X_{01}=E(X)_{01}$.
    Therefore $E^2\neq E$, so the doubled-tensor ZCL condition fails.
\end{proof}

\begin{remark}[Stabilized fixed points and the vertical algebra condition]%
    \label{rem:mpdo_algebra_tower_to_zcl_gap}
    \uses{def:mpdo_blocked_fixed_point_algebra_tower, def:mpo_is_zcl, thm:mpdo_algebra_tower_not_imply_zcl}
    Write $E=E_1$ for the one-site transfer map. Compatibility gives
    $\Fix((E^n)^\dagger)=\Fix(E^\dagger)$ for every $n>0$.
    Hence a nonzero matrix $X$ with $E^\dagger X=\lambda X$ and
    $\lambda^n=1$ lies in $\Fix(E^\dagger)$, so $(\lambda-1)X=0$ and
    $\lambda=1$ by
    Theorem~\ref{thm:mpdo_algebra_root_of_unity_eigenvalue}. Compatibility does
    not imply $E^2=E$, as
    Theorem~\ref{thm:mpdo_algebra_tower_not_imply_zcl} shows. More
    generally, on $\MN{D}$ with $D\ge2$, let $\Pi_{\mathrm{diag}}$ be the
    projection onto diagonal matrices. For $0<\varepsilon<1$, set
    $E(X)=\varepsilon X+(1-\varepsilon)\Pi_{\mathrm{diag}}(X)$, so that
    \begin{align}
        E^n(X)
         & = \Pi_{\mathrm{diag}}(X)
        + \varepsilon^n(X-\Pi_{\mathrm{diag}}(X)),
        \notag                           \\
        E^2-E
         & = (\varepsilon^2-\varepsilon)
        (\id-\Pi_{\mathrm{diag}}).
        \notag
    \end{align}
    The fixed-point algebra is diagonal for every $n>0$, but $E^2\ne E$.

    The coefficient argument in
    \cite[Appendix~C.4]{Cirac2016MPDO_arXiv} instead uses the vertical
    canonical form of an MPDO tensor $M$ in CF. Write its positive diagonal
    multiplicity matrices as $\mu_\alpha$ and its BNT as
    $\{M_\alpha\}_\alpha$. The operators $O_L(M_\alpha)$ are obtained by
    contracting $L$ tensors in the vertical direction, with the original
    physical and virtual roles exchanged. The required same-length product
    law is
    \begin{align}
        O_L(M_\alpha)O_L(M_\beta)
         & =
        \sum_\gamma
        \tr(\chi_{\alpha,\beta,\gamma}^{\,L})
        O_L(M_\gamma),
        \label{eq:mpdo_chi_product_law}
    \end{align}
    where the $\chi_{\alpha,\beta,\gamma}$ are positive diagonal matrices.
    With $m_\alpha=\tr(\mu_\alpha)$, the vector $m$ must be an idempotent
    for the multiplication induced by $c^{(1)}$, namely
    \begin{align}
        m_\gamma
         & =
        \sum_{\alpha,\beta}
        \tr(\chi_{\alpha,\beta,\gamma})
        m_\alpha m_\beta.
        \label{eq:mpdo_idempotent_trace_identity}
    \end{align}
    The proof of \cite[Theorem~4.14]{Cirac2016MPDO_arXiv} uses
    (\ref{eq:mpdo_chi_product_law}) and (\ref{eq:mpdo_idempotent_trace_identity}),
    positivity, and the simple-matrix lemma
    \cite[lines~1155--1163]{Cirac2016MPDO_arXiv} to identify the diagonal
    entries of the two-site multiplicity matrix $\nu_\gamma$ as a multiset:
    \begin{align}
        \{\nu_{\gamma,k}\}_k
         & =
        \{\mu_{\alpha,k_1}\mu_{\beta,k_2}
        \chi_{\alpha,\beta,\gamma,k_3}\}_{\alpha,\beta,k_1,k_2,k_3},
        \notag        \\
        \tr(\nu_\gamma)
         & =m_\gamma.
        \label{eq:mpdo_nu_chi_witness}
    \end{align}
    The reconstruction maps in \cite[lines~2065--2085]{Cirac2016MPDO_arXiv}
    then use
    \begin{align}
        \widetilde R_2\left(\bigoplus_\gamma X_\gamma\right)
         & = \bigoplus_\gamma
        \frac{\nu_\gamma}{m_\gamma}\otimes X_\gamma,
        \notag                \\
        \widetilde R_1\left(\bigoplus_\gamma X_\gamma\right)
         & = \bigoplus_\gamma
        \frac{\mu_\gamma}{m_\gamma}\otimes X_\gamma.
        \label{eq:mpdo_reconstruction_maps}
    \end{align}
    Thus $\tr(\nu_\gamma)=m_\gamma$ and $m_\gamma=\tr(\mu_\gamma)$
    normalize the factors $\nu_\gamma/m_\gamma$ and
    $\mu_\gamma/m_\gamma$, respectively, in
    $T=\widetilde R_2\widetilde T R_1$ and
    $S=\widetilde R_1\widetilde S R_2$. Here $R_1$ and $R_2$ compress to
    the canonical sectors and trace out their multiplicity factors, while
    $\widetilde T$ and $\widetilde S$ conjugate by the block unitaries
    relating the BNTs. These compositions are tpCPM on the retained spaces.
    When discarded zero-sector complements are present, extending the
    compressions by trace-and-prepare branches on the discarded input
    supports yields tpCPM on the full physical spaces. These additional
    branches vanish on the tensor letters, so the extended maps still satisfy
    $T[M_1(X)]=M_2(X)$ and $S[M_2(X)]=M_1(X)$, as required for an RFP.

    The scalar Newton--Girard identity relates traces of matrix powers to
    characteristic-polynomial coefficients; see
    the corresponding theorem in the companion quantum-channel volume~\cite[``Newton--Girard trace recursion'']{QICLean2026Blueprint}.
    Such scalar identities do not give the vertical canonical
    decomposition. The trace-power conclusion of
    \cite[Theorem~4.14(ii)]{Cirac2016MPDO_arXiv} concerns a common BNT and
    positive diagonal matrices satisfying
    (\ref{eq:mpdo_chi_product_law}) and
    (\ref{eq:mpdo_idempotent_trace_identity}). Relating these objects to the
    one-site and two-site vertical canonical forms is required for the maps
    in (\ref{eq:mpdo_reconstruction_maps}); it does not follow from a choice
    of bases in the support algebras. In particular,
    Theorem~\ref{thm:mpdo_algebra_tower_not_imply_zcl} rules out deriving
    idempotence of the doubled-tensor transfer map from compatibility alone.
    % Source-faithful completion requires the BNT-label hypotheses as a
    % separate condition; the support-algebra condition does not imply them.
\end{remark}

\subsection{Diagonal \texorpdfstring{$\chi$}{chi}-matrices and the trace-power formula}

In \cite[Theorem~4.14(ii)]{Cirac2016MPDO_arXiv}, the structure coefficients
of the length-$L$ operator algebra are traces of powers of positive diagonal
matrices $\chi_{\alpha,\beta,\gamma}$. In
Definition~\ref{def:mpdo_diagonal_chi_family}, arbitrary complex diagonal
entries are permitted to separate the scalar trace-power identity from
positivity and operator multiplication.

\begin{definition}[Diagonal \texorpdfstring{$\chi$}{chi} matrices]%
    \label{def:mpdo_diagonal_chi_family}
    \lean{MPOTensor.DiagonalChiFamily}
    \leanok
    Let $I$ be an index set. A family of diagonal complex matrices indexed
    by $(\alpha,\beta,\gamma)\in I^3$ consists of a size
    $r_{\alpha,\beta,\gamma}\in\N$ and entries
    $\chi_{\alpha,\beta,\gamma,k}\in\C$ for
    $0\le k<r_{\alpha,\beta,\gamma}$, defining
    \begin{align}
        \chi_{\alpha,\beta,\gamma}
         & =
        \diag\bigl(    (\chi_{\alpha,\beta,\gamma,k})_{0\le k<r_{\alpha,\beta,\gamma}}
        \bigr).
        \notag
    \end{align}
    A size of zero gives the empty diagonal matrix.
    The matrices $\chi_{\alpha,\beta,\gamma}$ in
    \cite[Theorem~4.14(ii)]{Cirac2016MPDO_arXiv} additionally have positive
    diagonal entries.
\end{definition}

\begin{theorem}[Trace of a power of a diagonal matrix]%
    \label{thm:mpdo_trace_matrix_pow}
    \lean{MPOTensor.DiagonalChiFamily.trace_matrix_pow}
    \leanok
    \uses{def:mpdo_diagonal_chi_family}
    For every triple $(\alpha, \beta, \gamma)$ and every $L \ge 0$,
    \begin{align}
        \tr(\chi_{\alpha,\beta,\gamma}^{L})
         & =
        \sum_{0\le k<r_{\alpha,\beta,\gamma}}
        \chi_{\alpha,\beta,\gamma,k}^{L},
        \notag
    \end{align}
    where $r_{\alpha,\beta,\gamma}$ is the size of
    $\chi_{\alpha,\beta,\gamma}$ and $\chi_{\alpha,\beta,\gamma,k}$ are its
    diagonal entries.
\end{theorem}
\begin{proof}\leanok
    Diagonal matrices raise to powers entrywise and the trace of a diagonal matrix
    is the sum of its entries.
\end{proof}

\begin{definition}[Trace-power form of structure coefficients]%
    \label{def:mpdo_has_chi_trace_power_form}
    \lean{MPOTensor.HasChiTracePowerForm}
    \leanok
    \uses{def:mpdo_diagonal_chi_family}
    Let $c^{(L)}_{\alpha,\beta,\gamma}$ be a family of complex coefficients
    and let $\chi_{\alpha,\beta,\gamma}$ be diagonal matrices on the same
    index set. The coefficients have trace-power form with respect to
    $\chi$ if, for every $L\ge0$ and every triple $(\alpha,\beta,\gamma)$,
    \begin{align}
        c^{(L)}_{\alpha,\beta,\gamma}
         & =
        \tr(\chi_{\alpha,\beta,\gamma}^{L})
        =
        \sum_{0\le k<r_{\alpha,\beta,\gamma}}
        \chi_{\alpha,\beta,\gamma,k}^{L}.
        \notag
    \end{align}
    Here $r_{\alpha,\beta,\gamma}$ is the size of
    $\chi_{\alpha,\beta,\gamma}$. This is the scalar identity appearing in
    \cite[Theorem~4.14(ii)]{Cirac2016MPDO_arXiv}. By itself, the definition
    neither requires positive diagonal entries nor asserts that the
    coefficients describe multiplication of vertical operators.
\end{definition}


\subsection{Structure coefficients and the idempotent condition}

\begin{definition}[Structure coefficients indexed by a BNT]%
    \label{def:mpdo_bnt_label_coefficient_family}
    \lean{MPOTensor.BNTLabelCoefficientFamily}
    \lean{MPOTensor.BNTLabelCoefficientFamily.ofChi}
    \lean{MPOTensor.BNTLabelCoefficientFamily.ofChi_coeff}
    \lean{MPOTensor.BNTLabelCoefficientFamily.ofChi_coeff_eq_trace_matrix_pow}
    \leanok
    Let $\Lambda$ be a fixed set of BNT labels. Consider complex coefficients
    $c^{(L)}_{\alpha,\beta,\gamma}$ for $L\in\N$ and
    $\alpha,\beta,\gamma\in\Lambda$. In
    \cite[Theorem~4.14(ii)]{Cirac2016MPDO_arXiv}, these are the structure
    coefficients for the products $O_L(M_\alpha)O_L(M_\beta)$ at the same
    length $L$, rather than for the blocked-basis multiplication
    $\mathcal A_n\times\mathcal A_n\to\mathcal A_{2n}$.
    A family of diagonal matrices $\chi_{\alpha,\beta,\gamma}$ determines
    coefficients
    $c^{(L)}_{\alpha,\beta,\gamma}
        = \tr(\chi_{\alpha,\beta,\gamma}^{L})
        = \sum_r\chi_{\alpha,\beta,\gamma,r}^{\,L}$.
    The positive diagonal matrices constructed in
    \cite[Appendix~C.4]{Cirac2016MPDO_arXiv} give coefficients of this form;
    the definition for a general diagonal family does not require positivity.
\end{definition}

\begin{definition}[Operators indexed by a BNT]%
    \label{def:mpdo_bnt_label_operator_family}
    \lean{MPOTensor.BNTLabelOperatorFamily}
    \leanok
    For each length $L$, let $O_L(M_\alpha)$ be operators indexed by
    the fixed BNT labels $\alpha$. Their ambient complex vector space,
    equipped with multiplication, may depend on $L$ and is not identified
    with a chosen blocked support algebra.
    Write $\widehat M_\alpha^{ab}=M_\alpha^{(a,b)}$ for the vertical reading.
    The vertical indices of the $L$ copies are contracted cyclically, while
    the horizontal legs remain open, as in
    \cite[lines~962--967]{Cirac2016MPDO_arXiv}.
    For $\mathbf a=(a_1,\ldots,a_L)$ and $\mathbf b=(b_1,\ldots,b_L)$,
    \begin{align}
        [O_L(M_\alpha)]_{\mathbf a,\mathbf b}
         & =
        \tr\!\left(\widehat M_\alpha^{a_1b_1}\cdots
        \widehat M_\alpha^{a_Lb_L}\right).
        \notag
    \end{align}
    Diagrammatically, with the cyclically contracted matrix index of each
    copy of $\widehat M_\alpha$ drawn horizontally,
    \begin{center}
        % Formula: [O_L(M_alpha)]_{a,b}=tr(\widehat M_alpha^{a_1b_1} ...
        % \widehat M_alpha^{a_Lb_L}), the display above; source
        % arXiv:1606.00608, lines 962--967 (figure MPDO_O_L_A_), drawn here
        % turned by a quarter turn, so the source's vertical closure is the
        % horizontal ring.
        % Ink-to-index map: each box is \widehat M_alpha; its west and east
        % legs are the matrix index of \widehat M_alpha^{ab}, which is the
        % vertical index of M; its north leg a_r and south leg b_r are the
        % horizontal indices of the r-th copy, which stay open.
        % Contraction set: the matrix indices, closed cyclically by the
        % periodic return (the trace); no a_r or b_r is summed.
        % Boundary signature: (virtual-west=0 | virtual-east=0 |
        % physical-up=3 | physical-down=3) for the schematic shown.
        \tenkzkernel
        \begin{tenkz}[rows={op}, cols=4, boundary=periodic]
            \tn[skin=mpo, size=l,
                ports={90:physical:$a_1$, 270:physical:$b_1$}]{\widehat M_\alpha} &
            \tn[skin=mpo, size=l,
                ports={90:physical:$a_2$, 270:physical:$b_2$}]{\widehat M_\alpha} &
            \tn[skin=dots, ports={180:virtual, 0:virtual}]{} &
            \tn[skin=mpo, size=l,
                ports={90:physical:$a_L$, 270:physical:$b_L$}]{\widehat M_\alpha}
        \end{tenkz}
    \end{center}
\end{definition}

\begin{definition}[Product of BNT operators at the same length]%
    \label{def:mpdo_bnt_label_same_length_product_form}
    \lean{MPOTensor.BNTLabelOperatorFamily.HasSameLengthProductForm}
    \lean{MPOTensor.BNTLabelOperatorFamily.HasSameLengthProductForm.eq_sum}
    \leanok
    \uses{def:mpdo_bnt_label_operator_family, def:mpdo_bnt_label_coefficient_family}
    For a finite set of BNT labels, the operators $O_L(M_\alpha)$ satisfy
    the product identity with coefficients $c^{(L)}_{\alpha,\beta,\gamma}$
    when, for every positive length $L$ and all labels $\alpha,\beta$,
    \begin{align}
        O_L(M_\alpha)O_L(M_\beta)
         & =
        \sum_\gamma
        c^{(L)}_{\alpha,\beta,\gamma} O_L(M_\gamma).
        \notag
    \end{align}
    In the diagram of Definition~\ref{def:mpdo_bnt_label_operator_family},
    the operator product joins the south legs of the first ring to the north
    legs of the second, and the identity reads
    \begin{tenkzequation}
        % Formula: O_L(M_alpha)O_L(M_beta)=sum_gamma c^{(L)}_{alpha,beta,gamma}
        % O_L(M_gamma), the display above; source arXiv:1606.00608,
        % Theorem 4.14(ii), lines 972--985 (eq:algebra).
        % Ink-to-index map: the upper ring is O_L(M_alpha) with open north
        % legs a_r, the lower ring is O_L(M_beta) with open south legs b_r;
        % the right ring is O_L(M_gamma).  Each ring closes the matrix
        % indices of its own boxes.
        % Contraction set: on the left, the vertical wire at each site is the
        % operator product, summing the south index of the alpha copy
        % against the north index of the beta copy; the two rings share no
        % other index.  On the right only the ring closure is summed.
        % Boundary signature: (virtual-west=0 | virtual-east=0 |
        % physical-up=3 | physical-down=3) on both sides.
        \tenkzkernel
        \begin{tenkz}[rows={op,op}, cols=4, boundary=periodic, bonds=none]
            \tn[skin=mpo, size=l, name=a1, ports={90:physical:$a_1$,
                        270:physical, 180:virtual, 0:virtual}]{\widehat M_\alpha} &
            \tn[skin=mpo, size=l, name=a2, ports={90:physical:$a_2$,
                        270:physical, 180:virtual, 0:virtual}]{\widehat M_\alpha} &
            \tn[skin=dots, name=ad, ports={180:virtual, 0:virtual}]{} &
            \tn[skin=mpo, size=l, name=aL, ports={90:physical:$a_L$,
                        270:physical, 180:virtual, 0:virtual}]{\widehat M_\alpha} \\
            \tn[skin=mpo, size=l, name=b1, ports={90:physical,
                        270:physical:$b_1$, 180:virtual, 0:virtual}]{\widehat M_\beta} &
            \tn[skin=mpo, size=l, name=b2, ports={90:physical,
                        270:physical:$b_2$, 180:virtual, 0:virtual}]{\widehat M_\beta} &
            \tn[skin=dots, name=bd, ports={180:virtual, 0:virtual}]{} &
            \tn[skin=mpo, size=l, name=bL, ports={90:physical,
                        270:physical:$b_L$, 180:virtual, 0:virtual}]{\widehat M_\beta}
            \tnwire{a1.0}{a2.180} \tnwire{a2.0}{ad.180} \tnwire{ad.0}{aL.180}
            \tnwire{b1.0}{b2.180} \tnwire{b2.0}{bd.180} \tnwire{bd.0}{bL.180}
            \tnwire{a1.270}{b1.90} \tnwire{a2.270}{b2.90}
            \tnwire{aL.270}{bL.90}
        \end{tenkz}
        \quad $=\displaystyle\sum_\gamma c^{(L)}_{\alpha,\beta,\gamma}$ \quad
        \begin{tenkz}[rows={op}, cols=4, boundary=periodic]
            \tn[skin=mpo, size=l,
                ports={90:physical:$a_1$, 270:physical:$b_1$}]{\widehat M_\gamma} &
            \tn[skin=mpo, size=l,
                ports={90:physical:$a_2$, 270:physical:$b_2$}]{\widehat M_\gamma} &
            \tn[skin=dots, ports={180:virtual, 0:virtual}]{} &
            \tn[skin=mpo, size=l,
                ports={90:physical:$a_L$, 270:physical:$b_L$}]{\widehat M_\gamma}
        \end{tenkz}
    \end{tenkzequation}
    This is the product identity appearing in
    \cite[Theorem~4.14(ii)]{Cirac2016MPDO_arXiv}, stated independently of
    the blocked-basis multiplication
    $\mathcal A_n\times\mathcal A_n\to\mathcal A_{2n}$.
\end{definition}

\begin{definition}[Scalars indexed by a BNT]%
    \label{def:mpdo_bnt_label_trace_scalar_family}
    \lean{MPOTensor.BNTLabelTraceScalarFamily}
    \leanok
    Let $m_\alpha\in\C$ be scalars indexed by the fixed BNT labels
    $\alpha$. For the positive diagonal matrices $\mu_\alpha$ in the
    vertical decomposition, these scalars are $m_\alpha=\tr(\mu_\alpha)$,
    as in the idempotent condition of
    \cite[Theorem~4.14(ii)]{Cirac2016MPDO_arXiv}.
\end{definition}

\begin{definition}[Idempotent for the multiplication induced by the coefficients]%
    \label{def:mpdo_bnt_label_idempotent_coefficient_form}
    \lean{MPOTensor.BNTLabelTraceScalarFamily.HasIdempotentCoefficientForm}
    \lean{MPOTensor.BNTLabelTraceScalarFamily.HasIdempotentCoefficientForm.eq_sum}
    \leanok
    \uses{def:mpdo_bnt_label_trace_scalar_family, def:mpdo_bnt_label_coefficient_family}
    For a finite set of BNT labels, the vector $m=(m_\alpha)_\alpha$ is an
    idempotent for the multiplication induced by $c^{(1)}$ when, for every
    label $\gamma$,
    \begin{align}
        m_\gamma
         & =
        \sum_{\alpha,\beta}
        c^{(1)}_{\alpha,\beta,\gamma}m_\alpha m_\beta.
        \notag
    \end{align}
    This is the idempotent condition in
    \cite[Theorem~4.14(ii)]{Cirac2016MPDO_arXiv}, stated separately from
    the trace-power formula.
\end{definition}

\begin{definition}[Assignment of BNT labels to blocked bases]%
    \label{def:mpdo_bnt_blocked_basis_label_assignment}
    \lean{MPOTensor.BNTBlockedBasisLabelAssignment}
    \lean{MPOTensor.BNTBlockedBasisLabelAssignment.blockedLabel}
    \leanok
    \uses{def:mpdo_blocked_structure_coefficients}
    An assignment of BNT labels consists, for every positive blocking length
    $n$, of a map $\sigma_n$ from the chosen basis indices of
    $\mathcal A_n$ and a map $\tau_n$ from the chosen basis indices of
    $\mathcal A_{2n}$ to the fixed set of BNT labels. The combined
    map $\sigma_n \sqcup \tau_n$ is defined on the disjoint union of these
    two index sets.
\end{definition}

\begin{definition}[Comparison of BNT structure coefficients with blocked coefficients]%
    \label{def:mpdo_bnt_blocked_basis_coefficient_comparison}
    \lean{MPOTensor.BNTBlockedBasisCoefficientComparison}
    \lean{MPOTensor.BNTBlockedBasisCoefficientComparison.sourceLabel}
    \lean{MPOTensor.BNTBlockedBasisCoefficientComparison.targetLabel}
    \lean{MPOTensor.BNTBlockedBasisCoefficientComparison.blocked_coeff_eq}
    \leanok
    \uses{def:mpdo_bnt_label_coefficient_family, def:mpdo_blocked_structure_coefficients, def:mpdo_bnt_blocked_basis_label_assignment}
    A comparison between the BNT structure coefficients and the chosen
    blocked-basis coefficients consists of an assignment of BNT labels
    together with the equality
    \begin{align}
        c^{(n)}_{i,j,k}
         & =
        c^{(n)}_{\sigma_n(i),\sigma_n(j),\tau_n(k)}
        \notag
    \end{align}
    for every $n>0$ and all basis indices $i,j$ of $\mathcal A_n$ and $k$
    of $\mathcal A_{2n}$.
    On the left, $c^{(n)}_{i,j,k}$ is the coefficient of the $k$-th basis
    element of $\mathcal A_{2n}$ in the blocked product. This comparison is
    an additional hypothesis, not the product identity for BNT operators
    at the same length.
\end{definition}

\begin{definition}[Trace-power formula at positive lengths]%
    \label{def:mpdo_bnt_label_positive_length_chi_trace_power_form}
    \lean{MPOTensor.BNTLabelCoefficientFamily.HasPositiveLengthChiTracePowerForm}
    \leanok
    \uses{def:mpdo_bnt_label_coefficient_family, def:mpdo_diagonal_chi_family}
    The coefficients $c^{(L)}_{\alpha,\beta,\gamma}$ satisfy the trace-power
    formula for a family of diagonal matrices $\chi_{\alpha,\beta,\gamma}$
    when, for every positive length $L$ and all labels,
    \begin{align}
        c^{(L)}_{\alpha,\beta,\gamma}
         & =
        \tr(\chi_{\alpha,\beta,\gamma}^{L}).
        \notag
    \end{align}
    The same matrix $\chi_{\alpha,\beta,\gamma}$ is used for all positive
    $L$. Its diagonal entries may be complex; positivity is not required
    here.
\end{definition}

\begin{lemma}[Structure coefficient as a trace power at positive length]%
    \label{lem:mpdo_bnt_label_chi_eq_trace_matrix_pow}
    \lean{MPOTensor.BNTLabelCoefficientFamily.HasPositiveLengthChiTracePowerForm.eq_trace_matrix_pow}
    \leanok
    \uses{def:mpdo_bnt_label_positive_length_chi_trace_power_form, thm:mpdo_trace_matrix_pow}
    Under the trace-power condition at positive lengths, for every $L>0$,
    $c^{(L)}_{\alpha,\beta,\gamma}
        = \tr(\chi_{\alpha,\beta,\gamma}^{L})$.
\end{lemma}
\begin{proof}\leanok
    Combining the defining positive-length trace-power identity with the trace
    formula for powers of a diagonal matrix gives the identity.
\end{proof}

\begin{lemma}[Coefficients defined by traces of diagonal matrix powers]%
    \label{lem:mpdo_bnt_label_of_chi_positive_length_trace_power}
    \lean{MPOTensor.BNTLabelCoefficientFamily.ofChi_hasPositiveLengthChiTracePowerForm}
    \leanok
    \uses{def:mpdo_bnt_label_coefficient_family, def:mpdo_bnt_label_positive_length_chi_trace_power_form}
    The coefficients defined by
    $c^{(L)}_{\alpha,\beta,\gamma}=\tr(\chi_{\alpha,\beta,\gamma}^{L})$
    for a family of diagonal matrices satisfy the trace-power condition at
    positive lengths for those matrices.
\end{lemma}
\begin{proof}\leanok
    This follows from the definition of the coefficients.
\end{proof}

\begin{definition}[Positive diagonal matrices giving the structure coefficients]%
    \label{def:mpdo_positive_bnt_label_chi_trace_power_form}
    \lean{MPOTensor.PositiveBNTLabelChiTracePowerForm}
    \lean{MPOTensor.PositiveBNTLabelChiTracePowerForm.ofChi}
    \leanok
    \uses{def:mpdo_bnt_label_positive_length_chi_trace_power_form, def:mpdo_diagonal_chi_family, lem:mpdo_bnt_label_of_chi_positive_length_trace_power}
    For given coefficients $c^{(L)}_{\alpha,\beta,\gamma}$, consider diagonal
    matrices $\chi_{\alpha,\beta,\gamma}$, independent of $L$, with strictly
    positive diagonal entries, such that
    $c^{(L)}_{\alpha,\beta,\gamma}=\tr(\chi_{\alpha,\beta,\gamma}^{L})$
    for every $L>0$ and all labels. This is the coefficient condition in
    \cite[Theorem~4.14(ii)]{Cirac2016MPDO_arXiv}.
    If the coefficients are defined by these traces, positivity of the
    diagonal entries suffices.
\end{definition}

\begin{lemma}[Trace powers of the positive diagonal matrices]%
    \label{lem:mpdo_positive_bnt_label_chi_eq_trace_matrix_pow}
    \lean{MPOTensor.PositiveBNTLabelChiTracePowerForm.eq_trace_pow}
    \leanok
    \uses{def:mpdo_positive_bnt_label_chi_trace_power_form, lem:mpdo_bnt_label_chi_eq_trace_matrix_pow}
    If the coefficients have a trace-power representation by positive
    diagonal matrices $\chi_{\alpha,\beta,\gamma}$ independent of $L$, then
    for every $L>0$,
    $c^{(L)}_{\alpha,\beta,\gamma}
        =\tr(\chi_{\alpha,\beta,\gamma}^{L})$.
\end{lemma}
\begin{proof}\leanok
    The trace reformulation of the trace-power identity at positive lengths
    gives the equality.
\end{proof}

\begin{lemma}[Product of BNT operators with trace-power coefficients]%
    \label{thm:mpdo_bnt_label_same_length_product_eq_sum_chi_trace_pow}
    \lean{MPOTensor.BNTLabelOperatorFamily.HasSameLengthProductForm.eq_sum_chi_trace_pow}
    \leanok
    \uses{def:mpdo_bnt_label_same_length_product_form, lem:mpdo_positive_bnt_label_chi_eq_trace_matrix_pow}
    Suppose the product identity for BNT operators at the same length holds
    and its coefficients have a trace-power representation by positive
    diagonal matrices $\chi_{\alpha,\beta,\gamma}$ independent of $L$.
    Then, for every $L>0$,
    \begin{align}
        O_L(M_\alpha)O_L(M_\beta)
         & =
        \sum_\gamma
        \tr(\chi_{\alpha,\beta,\gamma}^{\,L})
        O_L(M_\gamma).
        \notag
    \end{align}
\end{lemma}
\begin{proof}\leanok
    Substituting the positive-length trace-power formula into the same-length BNT
    product expansion yields the displayed sum.
\end{proof}

\begin{lemma}[Product of BNT operators for coefficients defined by traces]%
    \label{thm:mpdo_bnt_label_same_length_product_eq_sum_of_chi_trace_pow}
    \lean{MPOTensor.BNTLabelOperatorFamily.HasSameLengthProductForm.eq_sum_ofChi_trace_pow}
    \leanok
    \uses{def:mpdo_bnt_label_same_length_product_form, def:mpdo_bnt_label_coefficient_family}
    If the product identity for BNT operators at the same length holds with
    coefficients defined by traces of powers of diagonal matrices
    $\chi_{\alpha,\beta,\gamma}$, then for every $L>0$,
    \begin{align}
        O_L(M_\alpha)O_L(M_\beta)
         & =
        \sum_\gamma
        \tr(\chi_{\alpha,\beta,\gamma}^{\,L})
        O_L(M_\gamma).
        \notag
    \end{align}
\end{lemma}
\begin{proof}\leanok
    Substituting the defining trace-power formula for the coefficients into the
    same-length BNT product expansion yields the displayed sum.
\end{proof}

\begin{lemma}[Idempotent condition in terms of traces]%
    \label{thm:mpdo_bnt_label_idempotent_eq_sum_chi_trace}
    \lean{MPOTensor.BNTLabelTraceScalarFamily.HasIdempotentCoefficientForm.eq_sum_chi_trace}
    \leanok
    \uses{def:mpdo_bnt_label_idempotent_coefficient_form, lem:mpdo_positive_bnt_label_chi_eq_trace_matrix_pow}
    If $m$ is an idempotent for the multiplication induced by $c^{(1)}$ and
    the coefficients have a trace-power representation by positive diagonal
    matrices $\chi_{\alpha,\beta,\gamma}$ independent of $L$, then
    \begin{align}
        m_\gamma
         & =
        \sum_{\alpha,\beta}
        \tr(\chi_{\alpha,\beta,\gamma})
        m_\alpha m_\beta.
        \notag
    \end{align}
\end{lemma}
\begin{proof}\leanok
    Substituting the length-one trace-power formula into the idempotent
    coefficient identity gives the result.
\end{proof}

\begin{lemma}[Idempotent condition for coefficients defined by traces]%
    \label{thm:mpdo_bnt_label_idempotent_eq_sum_of_chi_trace}
    \lean{MPOTensor.BNTLabelTraceScalarFamily.HasIdempotentCoefficientForm.eq_sum_ofChi_trace}
    \leanok
    \uses{def:mpdo_bnt_label_idempotent_coefficient_form, def:mpdo_bnt_label_coefficient_family}
    If $m$ is an idempotent for the multiplication induced by $c^{(1)}$,
    where the coefficients are defined by traces of powers of diagonal
    matrices $\chi_{\alpha,\beta,\gamma}$, then
    \begin{align}
        m_\gamma
         & =
        \sum_{\alpha,\beta}
        \tr(\chi_{\alpha,\beta,\gamma})
        m_\alpha m_\beta.
        \notag
    \end{align}
\end{lemma}
\begin{proof}\leanok
    Substituting the defining length-one trace formula into the idempotent
    coefficient identity gives the result.
\end{proof}

\begin{lemma}[Blocked coefficients inherit BNT trace powers]%
    \label{lem:mpdo_bnt_blocked_basis_coefficient_eq_trace_pow}
    \lean{MPOTensor.BNTBlockedBasisCoefficientComparison.blocked_coeff_eq_trace_pow}
    \lean{MPOTensor.BNTBlockedBasisCoefficientComparison.%
        blocked_coeff_eq_ofChi_trace_pow}
    \lean{MPOTensor.BNTBlockedBasisCoefficientComparison.%
        blocked_coeff_eq_pulledBlockedChi_trace_pow}
    \leanok
    \uses{def:mpdo_bnt_blocked_basis_coefficient_comparison, def:mpdo_positive_bnt_label_chi_trace_power_form, lem:mpdo_positive_bnt_label_chi_eq_trace_matrix_pow}
    Suppose the chosen blocked-basis coefficients satisfy the comparison
    equality with the BNT structure coefficients, and the latter have a
    trace-power representation by positive diagonal matrices
    $\chi_{\alpha,\beta,\gamma}$ independent of length. Then, for every
    $n>0$ and every multiplication triple $(i,j,k)$,
    \begin{align}
        c^{(n)}_{i,j,k}
         & =
        \tr\!\left(\chi_{\sigma_n(i),\sigma_n(j),\tau_n(k)}^{\,n}\right).
        \notag
    \end{align}
\end{lemma}
\begin{proof}\leanok
    The comparison equality replaces the blocked-basis coefficient by the
    corresponding BNT structure coefficient, and the trace-power formula then
    gives the displayed trace.
    In the special case where the compared BNT coefficients are
    defined by traces of powers of $\chi$, the same conclusion
    follows directly from their definition.
\end{proof}

\begin{definition}[Diagonal matrices for blocked multiplication coefficients]%
    \label{def:mpdo_blocked_structure_chi_family}
    \lean{MPOTensor.AlgebraStructureData.BlockedStructureChiFamily}
    \leanok
    \uses{def:mpdo_blocked_structure_coefficients}
    For each positive blocking length $n$, choose a family of diagonal
    matrices indexed by triples in the disjoint union of the chosen basis
    indices of $\mathcal A_n$ and $\mathcal A_{2n}$. For a multiplication
    triple $(i,j,k)$, where $i,j$ label basis elements of $\mathcal A_n$ and
    $k$ labels a basis element of $\mathcal A_{2n}$, write its size as
    $r_{n,i,j,k}\in\N$ and its complex diagonal entries as
    $\chi_{n,i,j,k,s}$, with $0\le s<r_{n,i,j,k}$. The corresponding matrix is
    \begin{align}
        \chi_{n,i,j,k}
         & =
        \diag(\chi_{n,i,j,k,s})_{s=0}^{r_{n,i,j,k}-1}.
        \notag
    \end{align}
    Unlike the matrices indexed by fixed BNT labels in
    \cite[Theorem~4.14(ii)]{Cirac2016MPDO_arXiv}, the size and diagonal
    entries may depend on the blocking length and on the chosen basis indices.
\end{definition}
% The uniformity gap is documented in
% docs/paper-gaps/cpgsv17_blocked_chi_uniformity.tex.

\begin{lemma}[Trace of a power of a blocked diagonal matrix]%
    \label{lem:mpdo_blocked_structure_chi_trace_matrix_pow}
    \lean{MPOTensor.AlgebraStructureData.BlockedStructureChiFamily.trace_matrix_pow}
    \leanok
    \uses{def:mpdo_blocked_structure_chi_family}
    For each positive blocking length $n$, each multiplication triple $(i,j,k)$,
    and every exponent $L\in\N$,
    $\tr(\chi_{n,i,j,k}^{\,L})
        = \sum_{s=0}^{r_{n,i,j,k}-1} \chi_{n,i,j,k,s}^{\,L}$.
\end{lemma}
\begin{proof}\leanok
    Diagonal matrices raise to powers entrywise and the trace of a diagonal
    matrix is the sum of its entries.
\end{proof}

\begin{definition}[Trace-power form for blocked multiplication coefficients]%
    \label{def:mpdo_has_blocked_structure_chi_trace_power_form}
    \lean{MPOTensor.AlgebraStructureData.HasBlockedStructureChiTracePowerForm}
    \leanok
    \uses{def:mpdo_blocked_structure_coefficients, def:mpdo_blocked_structure_chi_family}
    The blocked multiplication coefficients are in trace-power form when, for
    every positive blocking length $n$ and every multiplication triple $(i,j,k)$,
    the coefficient $c^{(n)}_{i,j,k}$ equals
    $\sum_{s=0}^{r_{n,i,j,k}-1} \chi_{n,i,j,k,s}^{\,n}$.
    Equivalently, by
    Lemma~\ref{lem:mpdo_blocked_structure_chi_trace_matrix_pow}, it is the
    trace of the $n$-th power of the corresponding diagonal matrix. The exponent
    $n$ is the blocking length of the two factors in $\mathcal A_n$; the
    product is expanded in the basis of $\mathcal A_{2n}$.
\end{definition}

\begin{lemma}[Blocked coefficient as a trace power]%
    \label{lem:mpdo_blocked_structure_chi_eq_trace_matrix_pow}
    \lean{MPOTensor.AlgebraStructureData.HasBlockedStructureChiTracePowerForm.eq_trace_matrix_pow}
    \lean{MPOTensor.AlgebraStructureData.PositiveBlockedStructureChiTracePowerForm.eq_trace_pow}
    \leanok
    \uses{def:mpdo_has_blocked_structure_chi_trace_power_form, lem:mpdo_blocked_structure_chi_trace_matrix_pow}
    Under trace-power form for the blocked multiplication coefficients,
    for every positive blocking length $n$,
    $c^{(n)}_{i,j,k}=\tr(\chi_{n,i,j,k}^{\,n})$.
\end{lemma}
\begin{proof}\leanok
    Combining the defining trace-power identity with the trace formula for
    powers of a diagonal matrix gives the equality.
\end{proof}

\begin{definition}[Positive diagonal matrices giving the blocked coefficients]%
    \label{def:mpdo_positive_blocked_structure_chi_trace_power_form}
    \lean{MPOTensor.AlgebraStructureData.PositiveBlockedStructureChiTracePowerForm}
    \leanok
    \uses{def:mpdo_has_blocked_structure_chi_trace_power_form, def:mpdo_blocked_structure_chi_family}
    Consider a family of blocked diagonal matrices with strictly positive
    diagonal entries satisfying
    $c^{(n)}_{i,j,k}=\tr(\chi_{n,i,j,k}^{\,n})$ for every $n>0$ and every
    multiplication triple $(i,j,k)$. This is a blocked-basis analogue of the
    positive diagonal matrices in
    \cite[Theorem~4.14(ii)]{Cirac2016MPDO_arXiv}, without requiring a single
    family indexed only by the BNT labels.
\end{definition}
% The uniformity gap is documented in
% docs/paper-gaps/cpgsv17_blocked_chi_uniformity.tex.

\begin{lemma}[Positive BNT diagonal matrices give positive blocked diagonal matrices]%
    \label{thm:mpdo_bnt_blocked_basis_to_positive_blocked_chi_trace_power_form}
    \lean{MPOTensor.BNTBlockedBasisCoefficientComparison.blockedLabel}
    \lean{MPOTensor.BNTBlockedBasisCoefficientComparison.pulledBlockedChiFamily}
    \lean{MPOTensor.BNTBlockedBasisCoefficientComparison.pulledBlockedChiFamily_toDiagonal_of_pos}
    \lean{MPOTensor.BNTBlockedBasisCoefficientComparison.%
        pulledBlockedChiFamily_toDiagonal_posEntries}
    \lean{MPOTensor.BNTBlockedBasisCoefficientComparison.pulledBlockedChi_tracePowerCoeff_of_pos}
    \lean{MPOTensor.BNTBlockedBasisCoefficientComparison.pulledBlockedChi_dim_of_pos}
    \lean{MPOTensor.BNTBlockedBasisCoefficientComparison.pulledBlockedChi_trace_matrix_pow_of_pos}
    \lean{MPOTensor.BNTBlockedBasisCoefficientComparison.%
        toPositiveBlockedStructureChiTracePowerForm}
    \leanok
    \uses{def:mpdo_bnt_blocked_basis_coefficient_comparison, def:mpdo_positive_bnt_label_chi_trace_power_form, def:mpdo_positive_blocked_structure_chi_trace_power_form}
    Suppose the chosen blocked-basis coefficients satisfy the comparison
    equality with the BNT structure coefficients, and the latter have a
    trace-power representation by positive diagonal matrices
    $\chi_{\alpha,\beta,\gamma}$ independent of length. Then there are
    positive blocked diagonal matrices giving the blocked coefficients by
    the trace-power formula.
    They are obtained by reindexing the BNT diagonal matrices along
    $\sigma_n\sqcup\tau_n$; on multiplication triples this gives
    $\chi_{n,i,j,k}=\chi_{\sigma_n(i),\sigma_n(j),\tau_n(k)}$.
\end{lemma}
\begin{proof}\leanok
    Reindexing the positive diagonal matrices along the blocked-basis
    label map preserves positivity, and the trace-power
    identity follows by substituting the blocked-basis comparison into the
    BNT trace-power formula.
\end{proof}

\begin{lemma}[Constant power sums have entries zero or one]%
    \label{lem:constant_power_sum_dichotomy}
    \lean{le_one_of_forall_sum_pow_eq}
    \lean{eq_zero_or_eq_one_of_forall_sum_pow_eq}
    \lean{Complex.eq_zero_or_eq_one_of_forall_sum_pow_eq}
    \leanok
    Let $x_0,\ldots,x_{n-1}$ be non-negative real numbers whose power sums
    $\sum_k x_k^L$ take the same value at every integer exponent $L \ge 1$.
    Then every $x_k$ equals $0$ or $1$. The same holds for a finite family of
    complex numbers whose imaginary parts vanish and whose real parts are
    non-negative.
\end{lemma}
\begin{proof}\leanok
    If some entry exceeded one, its powers would grow without bound while
    all other summands stay non-negative, contradicting constancy of the
    power sums; hence every entry lies in $[0,1]$. Comparing the exponents
    one and two gives $\sum_k x_k(1-x_k) = 0$ with non-negative summands, so
    each summand vanishes.
\end{proof}

\begin{definition}[Length-independent coefficients]%
    \label{def:mpdo_bnt_label_length_independent}
    \lean{MPOTensor.BNTLabelCoefficientFamily.LengthIndependent}
    \lean{MPOTensor.BNTLabelCoefficientFamily.LengthIndependent.coeff_eq}
    \leanok
    \uses{def:mpdo_bnt_label_coefficient_family}
    The structure coefficients $c^{(L)}_{\alpha,\beta,\gamma}$ are
    \emph{length independent} when
    $c^{(L)}_{\alpha,\beta,\gamma} = c^{(1)}_{\alpha,\beta,\gamma}$ for
    every positive length $L$ and all labels; equivalently, the
    coefficients agree at any two positive lengths. This is the case
    singled out in the discussion following
    \cite[Theorem~4.14]{Cirac2016MPDO_arXiv}.
\end{definition}

\begin{lemma}[Length-independent coefficients are natural numbers]%
    \label{thm:mpdo_bnt_label_length_independent_natural}
    \lean{MPOTensor.DiagonalChiFamily.%
        entry_eq_one_of_posEntries_of_forall_tracePowerCoeff_eq}
    \lean{MPOTensor.DiagonalChiFamily.matrix_eq_one_of_forall_entry_eq_one}
    \lean{MPOTensor.BNTLabelCoefficientFamily.HasPositiveLengthChiTracePowerForm.%
        entry_eq_one_of_lengthIndependent}
    \lean{MPOTensor.BNTLabelCoefficientFamily.HasPositiveLengthChiTracePowerForm.%
        coeff_eq_dim_of_lengthIndependent}
    \lean{MPOTensor.BNTLabelCoefficientFamily.HasPositiveLengthChiTracePowerForm.%
        exists_natCast_coeff_of_lengthIndependent}
    \lean{MPOTensor.PositiveBNTLabelChiTracePowerForm.entry_eq_one_of_lengthIndependent}
    \lean{MPOTensor.PositiveBNTLabelChiTracePowerForm.coeff_eq_dim_of_lengthIndependent}
    \leanok
    \uses{def:mpdo_bnt_label_length_independent, def:mpdo_bnt_label_positive_length_chi_trace_power_form, def:mpdo_positive_bnt_label_chi_trace_power_form, lem:constant_power_sum_dichotomy}
    Suppose the BNT structure coefficients have the trace-power form
    $c^{(L)}_{\alpha,\beta,\gamma}
        = \tr(\chi_{\alpha,\beta,\gamma}^{\,L})$
    with positive diagonal matrices $\chi_{\alpha,\beta,\gamma}$, and are
    length independent. Then every diagonal entry
    $\chi_{\alpha,\beta,\gamma,r}$ lies in $\{0,1\}$, as stated in
    \cite[Section~4.5]{Cirac2016MPDO_arXiv}; by positivity every entry
    equals one. Consequently, at every positive length,
    $c^{(L)}_{\alpha,\beta,\gamma}
        = r_{\alpha,\beta,\gamma} \in \N$, where
    $r_{\alpha,\beta,\gamma}$ is the size of the diagonal matrix.
    The structure coefficient is this multiplicity, not necessarily one.
\end{lemma}
\begin{proof}\leanok
    For fixed labels, the trace-power form turns length independence into
    constancy of the power sums of the diagonal entries of
    $\chi_{\alpha,\beta,\gamma}$, so by
    Lemma~\ref{lem:constant_power_sum_dichotomy} every entry equals $0$ or
    $1$, and positivity excludes $0$. The trace of
    $\chi_{\alpha,\beta,\gamma}^{\,L}$ is then the number of diagonal
    entries.
\end{proof}

\begin{lemma}[Unit entries give length independence]%
    \label{thm:mpdo_bnt_label_length_independent_of_unit_entries}
    \lean{MPOTensor.DiagonalChiFamily.tracePowerCoeff_eq_dim_of_forall_entry_eq_one}
    \lean{MPOTensor.BNTLabelCoefficientFamily.HasPositiveLengthChiTracePowerForm.%
        lengthIndependent_of_forall_entry_eq_one}
    \leanok
    \uses{def:mpdo_bnt_label_length_independent, def:mpdo_bnt_label_positive_length_chi_trace_power_form}
    If the BNT structure coefficients have a trace-power representation by
    diagonal matrices all of whose diagonal entries equal one, then for
    every $L>0$, $c^{(L)}_{\alpha,\beta,\gamma}
        = r_{\alpha,\beta,\gamma}$, so the coefficients are length
    independent.
\end{lemma}
\begin{proof}\leanok
    The power sums of a family of ones do not depend on the exponent.
\end{proof}

\begin{lemma}[Three consecutive positive power sums]%
    \label{thm:positive_power_second_difference}
    \lean{Real.sum_pow_secondDifference}
    \lean{Real.eq_one_of_sum_pow_secondDifference_eq_zero}
    \lean{Complex.eq_one_of_sum_pow_secondDifference_eq_zero}
    \lean{MPOTensor.DiagonalChiFamily.entry_eq_one_of_tracePowerCoeff_eq_at_three_lengths}
    \leanok
    Let $x_0,\ldots,x_{r-1}>0$ and $p_n=\sum_k x_k^n$.
    For every integer $n\geq0$,
    \begin{align}
        p_{n+2}-2p_{n+1}+p_n
        =\sum_k x_k^n(x_k-1)^2.
        \label{eq:mpdo_positive_second_difference}
    \end{align}
    The second difference vanishes if and only if every $x_k$ equals
    one. In particular, equality of three consecutive trace powers of
    a positive diagonal matrix forces that matrix to be the identity.
\end{lemma}
\begin{proof}\leanok
    Expanding the squares gives~(\ref{eq:mpdo_positive_second_difference}).
    Every summand is nonnegative, and every factor $x_k^n$ is strictly
    positive. Thus the sum vanishes exactly when all $x_k-1$ vanish.
\end{proof}

\begin{lemma}[Finite criterion for length independence]%
    \label{thm:mpdo_finite_length_independence}
    \lean{MPOTensor.BNTLabelCoefficientFamily.HasPositiveLengthChiTracePowerForm.%
        lengthIndependent_iff_coeff_two_three}
    \leanok
    \uses{def:mpdo_bnt_label_positive_length_chi_trace_power_form, thm:positive_power_second_difference, thm:mpdo_bnt_label_length_independent_of_unit_entries}
    Suppose the structure coefficients admit a trace-power presentation
    by positive diagonal matrices. They are length independent if and
    only if, for every triple of labels,
    $c^{(2)}_{\alpha,\beta,\gamma}=c^{(1)}_{\alpha,\beta,\gamma}$ and
    $c^{(3)}_{\alpha,\beta,\gamma}=c^{(1)}_{\alpha,\beta,\gamma}$.
\end{lemma}
\begin{proof}\leanok
    Length independence implies both equalities. Conversely, applying
    Lemma~\ref{thm:positive_power_second_difference} with $n=1$ to
    each diagonal matrix shows that all its entries equal one, so
    Lemma~\ref{thm:mpdo_bnt_label_length_independent_of_unit_entries}
    gives length independence.
\end{proof}

\input{chapter/ch21_mpdo_rfp_bnt_theorem_data_algebra_extensions.tex}%
\input{chapter/ch21_mpdo_rfp_bnt_theorem_data_singleton.tex}%
\input{chapter/ch21_mpdo_rfp_bnt_theorem_data_multiplicity_gauges.tex}%
\input{chapter/ch21_mpdo_rfp_bnt_theorem_data_idempotence_comparison.tex}%

\begin{definition}[BNT operators and structure coefficients]%
    \label{def:mpdo_bnt_label_theorem_witness}
    \lean{MPOTensor.BNTLabelTheoremWitness}
    \lean{MPOTensor.HasBNTLabelTheoremWitness}
    \lean{MPOTensor.BNTLabelTheoremWitness.toTheoremData}
    \lean{MPOTensor.BNTLabelTheoremWitness.coeffs}
    \lean{MPOTensor.BNTLabelTheoremWitness.operators}
    \lean{MPOTensor.BNTLabelTheoremWitness.traceScalars}
    \lean{MPOTensor.BNTLabelTheoremWitness.positiveChi}
    \lean{MPOTensor.BNTLabelTheoremWitness.blockedComparison}
    \lean{MPOTensor.BNTLabelTheoremWitness.same_length_product_form}
    \lean{MPOTensor.BNTLabelTheoremWitness.idempotent_coefficient_form}
    \lean{MPOTensor.BNTLabelTheoremWitness.same_length_product_form_ofChi}
    \lean{MPOTensor.BNTLabelTheoremWitness.idempotent_coefficient_form_ofChi}
    \lean{MPOTensor.BNTLabelTheoremWitness.same_length_product_eq_sum_ofChi}
    \lean{MPOTensor.BNTLabelTheoremWitness.idempotent_eq_sum_ofChi}
    \lean{MPOTensor.BNTLabelTheoremWitness.positive_chi_pos_entries}
    \lean{MPOTensor.BNTLabelTheoremWitness.positive_chi_trace_power}
    \lean{MPOTensor.BNTLabelTheoremWitness.labelAssignment}
    \lean{MPOTensor.BNTLabelTheoremWitness.sourceLabel}
    \lean{MPOTensor.BNTLabelTheoremWitness.targetLabel}
    \lean{MPOTensor.BNTLabelTheoremWitness.positiveBlockedChi}
    \lean{MPOTensor.BNTLabelTheoremWitness.ofChi}
    \lean{MPOTensor.BNTLabelTheoremWitness.blocked_coeff_eq}
    \lean{MPOTensor.BNTLabelTheoremWitness.blockedComparison_ofChi}
    \lean{MPOTensor.BNTLabelTheoremWitness.blocked_coeff_eq_ofChi}
    \leanok
    \uses{def:mpdo_bnt_label_theorem_data}
    Fix the algebra-structure data. Consider a finite set of labels $\alpha$,
    complex vector spaces equipped with a product at each length $L$, and
    operators $O_L(M_\alpha)$ in those spaces, together with scalars
    $m_\alpha$, structure coefficients $c^{(L)}_{\alpha,\beta,\gamma}$,
    and length-independent diagonal matrices
    $\chi_{\alpha,\beta,\gamma}$ with positive diagonal entries. Require
    the product, idempotent, trace-power, and blocked-basis comparison
    conditions of Definition~\ref{def:mpdo_bnt_label_theorem_data}.

    In a vertical basis-of-normal-tensors decomposition, $O_L(M_\alpha)$ denotes the operator
    obtained by contracting $L$ copies of $M_\alpha$, and
    $m_\alpha=\tr(\mu_\alpha)$ for the positive diagonal multiplicity
    matrix $\mu_\alpha$. The abstract conditions here do not themselves
    construct such a tensor decomposition.
    % The converse direction of CPSV16 Theorem 4.14 compares the algebra-side
    % and blocked physical sectors through marked/common-target tensors, proves
    % X_\gamma^\dagger X_\gamma=\omega_\gamma I with \omega_\gamma>0, normalizes
    % X_\gamma to U_\gamma=\omega_\gamma^{-1/2}X_\gamma, and constructs the
    % physical tpCP maps T=\widetilde R_2\widetilde T R_1 and
    % S=\widetilde R_1\widetilde S R_2. The mixed-prefix comparison supplies
    % the common target under the standing canonical-form and MPDO hypotheses.
    Existence means that these objects can be chosen simultaneously.
    The blocked-basis comparisons are additional conditions relating them
    to the fixed algebra-structure data, not clauses of the basis-of-normal-tensors
    algebra characterization.

    One may instead start with the positive diagonal matrices and define
    $c^{(L)}_{\alpha,\beta,\gamma}
        =\tr(\chi_{\alpha,\beta,\gamma}^{\,L})$,
    provided the product, idempotent, and blocked-basis comparison
    conditions hold. For any choice satisfying the conditions above,
    replacing the coefficients by these traces at positive lengths
    preserves all three conditions.
\end{definition}

\begin{theorem}[Product and idempotent equations for BNT operators]%
    \label{thm:mpdo_has_bnt_label_theorem_witness_source_equations}
    \lean{MPOTensor.HasBNTLabelTheoremWitness.exists_source_chi_trace_equations}
    \lean{MPOTensor.BNTLabelTheoremWitness.coeff_eq_trace_pow}
    \lean{MPOTensor.BNTLabelTheoremWitness.coeff_eq_ofChi_coeff}
    \lean{MPOTensor.BNTLabelTheoremWitness.chi_entry_pos}
    \lean{MPOTensor.BNTLabelTheoremWitness.same_length_product_eq_sum}
    \lean{MPOTensor.BNTLabelTheoremWitness.idempotent_eq_sum}
    \lean{MPOTensor.BNTLabelTheoremWitness.same_length_product_eq_sum_chi_trace_pow}
    \lean{MPOTensor.BNTLabelTheoremWitness.idempotent_eq_sum_chi_trace}
    \leanok
    \uses{def:mpdo_bnt_label_theorem_witness}
    Suppose the objects in
    Definition~\ref{def:mpdo_bnt_label_theorem_witness} exist.
    For any such choice, the product law holds for every positive length
    $L$ and all labels $\alpha,\beta$, and the scalar identity holds for
    every label $\gamma$:
    \begin{align}
        O_L(M_\alpha)O_L(M_\beta)
         & = \sum_\gamma c^{(L)}_{\alpha,\beta,\gamma} O_L(M_\gamma),
        \label{eq:mpdo_bnt_same_len_product}                                    \\
        m_\gamma
         & = \sum_{\alpha,\beta} c^{(1)}_{\alpha,\beta,\gamma}m_\alpha m_\beta.
        \label{eq:mpdo_bnt_idempotent_law}
    \end{align}
    Thus the vector $m=(m_\alpha)_\alpha$ is an idempotent for the
    multiplication induced by $c^{(1)}$.
    The diagonal matrices $\chi_{\alpha,\beta,\gamma}$ have positive
    diagonal entries and satisfy
    $c^{(L)}_{\alpha,\beta,\gamma}
        = \tr(\chi_{\alpha,\beta,\gamma}^{\,L})$
    for every $L>0$.
    Substituting this identity into~(\ref{eq:mpdo_bnt_same_len_product})
    and~(\ref{eq:mpdo_bnt_idempotent_law}) gives
    \begin{align}
        O_L(M_\alpha)O_L(M_\beta)
         & = \sum_\gamma \tr(\chi_{\alpha,\beta,\gamma}^{\,L}) O_L(M_\gamma),
        \label{eq:mpdo_bnt_same_len_product_chi}                                   \\
        m_\gamma
         & = \sum_{\alpha,\beta} \tr(\chi_{\alpha,\beta,\gamma}) m_\alpha m_\beta.
        \label{eq:mpdo_bnt_idempotent_law_chi}
    \end{align}
    At positive lengths, the coefficients therefore agree with those
    defined directly by the traces of powers of $\chi$.

    Let $\sigma_n$ and $\tau_n$ be the source and target label maps in
    the additional blocked-basis comparison at a positive length $n$.
    For source indices $i,j$ at length $n$ and a target index $k$ at
    length $2n$, the blocked-basis coefficients satisfy
    \begin{align}
        c^{(n)}_{i,j,k}
         & = c^{(n)}_{\sigma_n(i),\sigma_n(j),\tau_n(k)}
        = \tr\!\left(\chi_{\sigma_n(i),\sigma_n(j),\tau_n(k)}^{\,n}\right).
        \notag
    \end{align}
\end{theorem}
\begin{proof}\leanok
    For the assumed operators, scalars, coefficients, and diagonal matrices,
    the product and idempotent identities give
    (\ref{eq:mpdo_bnt_same_len_product})
    and~(\ref{eq:mpdo_bnt_idempotent_law}), and substituting the trace-power
    identity yields
    (\ref{eq:mpdo_bnt_same_len_product_chi})
    and~(\ref{eq:mpdo_bnt_idempotent_law_chi}).
    The blocked-basis comparison gives the final coefficient formula.
\end{proof}

\begin{theorem}[Positive diagonal matrices for blocked coefficients]%
    \label{thm:mpdo_bnt_label_theorem_witness_to_positive_blocked_chi_trace_power_form}
    \lean{MPOTensor.BNTLabelTheoremWitness.toPositiveBlockedStructureChiTracePowerForm}
    \lean{MPOTensor.BNTLabelTheoremWitness.positiveBlockedChi_toDiagonal_of_pos}
    \lean{MPOTensor.BNTLabelTheoremWitness.positiveBlockedChi_tracePowerCoeff_of_pos}
    \lean{MPOTensor.BNTLabelTheoremWitness.positiveBlockedChi_dim_of_pos}
    \lean{MPOTensor.BNTLabelTheoremWitness.positiveBlockedChi_trace_matrix_pow_of_pos}
    \lean{MPOTensor.BNTLabelTheoremWitness.positiveBlockedChi_tracePower}
    \lean{MPOTensor.BNTLabelTheoremWitness.positiveBlockedChi_posEntries}
    \lean{MPOTensor.BNTLabelTheoremWitness.positiveBlockedChi_entry_pos}
    \lean{MPOTensor.HasBNTLabelTheoremWitness.positive_blocked_chi_witness}
    \lean{MPOTensor.HasBNTLabelTheoremWitness.exists_blocked_chi_trace_power_form}
    \lean{MPOTensor.HasBNTLabelTheoremWitness.exists_blocked_coeff_eq_trace_pow}
    \leanok
    \uses{def:mpdo_bnt_label_theorem_witness, thm:mpdo_bnt_label_theorem_data_to_positive_blocked_chi_trace_power_form}
    Suppose the objects in
    Definition~\ref{def:mpdo_bnt_label_theorem_witness} exist for the
    chosen algebra-structure data. Then there are diagonal matrices
    $\chi_{n,i,j,k}$ with positive diagonal entries such that, for every
    positive blocked length $n$, source indices $i,j$ at length $n$,
    and target index $k$ at length $2n$,
    \begin{align}
        \chi_{n,i,j,k}
         & = \chi_{\sigma_n(i),\sigma_n(j),\tau_n(k)},
        \notag                                         \\
        c^{(n)}_{i,j,k}
         & = \tr(\chi_{n,i,j,k}^{\,n}).
        \notag
    \end{align}
    These matrices are obtained by selecting the BNT matrices
    $\chi_{\alpha,\beta,\gamma}$ along the source and target comparison
    maps $\sigma_n,\tau_n$. Their dimensions, diagonal entries, and
    traces of powers agree with those of the selected matrices.
    The same construction proves existence when the BNT operators,
    scalars, and matrices are assumed to exist without a fixed choice.
\end{theorem}
\begin{proof}\leanok
    Applying the blocked-basis trace-power consequence to the chosen
    operators, scalars, coefficients, and diagonal matrices yields the
    positive diagonal matrices $\chi_{n,i,j,k}$.
\end{proof}

\begin{corollary}[Blocked coefficients as traces of powers]%
    \label{cor:mpdo_bnt_label_theorem_witness_blocked_coeff_eq_positive_blocked_chi_trace_pow}
    \lean{MPOTensor.BNTLabelTheoremWitness.blocked_coeff_eq_trace_pow}
    \lean{MPOTensor.BNTLabelTheoremWitness.blocked_coeff_eq_positiveBlockedChi_trace_pow}
    \leanok
    \uses{def:mpdo_bnt_label_theorem_witness, thm:mpdo_bnt_label_theorem_witness_to_positive_blocked_chi_trace_power_form, lem:mpdo_blocked_structure_chi_eq_trace_matrix_pow}
    Assume the conditions of
    Definition~\ref{def:mpdo_bnt_label_theorem_witness}.
    For every positive blocked length $n$, source indices $i,j$ at
    length $n$, and target index $k$ at length $2n$, the blocked-basis
    coefficients are traces of powers of the BNT matrices selected by
    the source and target comparison maps:
    \begin{align}
        c^{(n)}_{i,j,k}
         & = \tr\!\left(\chi_{\sigma_n(i),\sigma_n(j),\tau_n(k)}^{\,n}\right)
        = \tr(\chi_{n,i,j,k}^{\,n}).
        \notag
    \end{align}
\end{corollary}
\begin{proof}\leanok
    The trace-power identity for the positive diagonal matrices
    $\chi_{n,i,j,k}$ obtained from the assumed BNT matrices and
    blocked-basis comparison gives the displayed equality.
\end{proof}

\subsection{Direct-sum sector conjugations}

\begin{definition}[Direct-sum unitary conjugations]
    \label{def:mpdo_bnt_direct_sum_unitary_conjugations}
    \lean{MPOTensor.BNTAlgebraTensorClause.TwoSiteExactSectorGauge.%
        RelabeledTwoSiteSectorAlgebra}
    \lean{MPOTensor.BNTAlgebraTensorClause.TwoSiteExactSectorGauge.%
        UnitarySectorConjugacy}
    \lean{MPOTensor.BNTAlgebraTensorClause.TwoSiteExactSectorGauge.%
        UnitarySectorConjugacy.sectorLinearEquiv}
    \lean{MPOTensor.BNTAlgebraTensorClause.TwoSiteExactSectorGauge.%
        UnitarySectorConjugacy.directSumUnitaryT}
    \lean{MPOTensor.BNTAlgebraTensorClause.TwoSiteExactSectorGauge.%
        UnitarySectorConjugacy.directSumUnitaryS}
    \lean{MPOTensor.BNTAlgebraTensorClause.TwoSiteExactSectorGauge.%
        UnitarySectorConjugacy.directSumUnitaryT_apply}
    \lean{MPOTensor.BNTAlgebraTensorClause.TwoSiteExactSectorGauge.%
        UnitarySectorConjugacy.directSumUnitaryS_apply}
    \leanok
    \uses{def:mpdo_bnt_two_site_exact_sector_gauge}
    Work in the setting of
    Definition~\ref{def:mpdo_bnt_two_site_exact_sector_gauge}.  Let
    \begin{align}
        \mathcal A_1
         & =\bigoplus_\gamma \MN{d_\gamma},
         &
        \mathcal A_2^\sigma
         & =\bigoplus_\gamma \MN{\widetilde d_{\sigma(\gamma)}}
    \end{align}
    be the one-site sector algebra and the two-site sector algebra relabelled
    by the given sector matching.  Suppose that
    $d_\gamma=\widetilde d_{\sigma(\gamma)}$ and that there is a unitary
    $U_\gamma$ on each matched sector such that
    \begin{align}
        A_{\sigma(\gamma)}^i
         & =U_\gamma M_\gamma^i U_\gamma^*
        \label{eq:mpdo_bnt_sector_unitary_conjugacy}
    \end{align}
    for every physical index $i$.  Denote by $\iota_\gamma$ the coordinate
    identification induced by the equality of the two bond dimensions.
    Define mutually inverse linear maps
    \begin{align}
        \widetilde T\colon \mathcal A_1        & \longrightarrow\mathcal A_2^\sigma,
                                               &
        [\widetilde T(X)]_\gamma
                                               & =U_\gamma\iota_\gamma(X_\gamma)U_\gamma^*,
        \label{eq:mpdo_bnt_direct_sum_unitary_T}
        \\
        \widetilde S\colon \mathcal A_2^\sigma & \longrightarrow\mathcal A_1,
                                               &
        [\widetilde S(Y)]_\gamma
                                               & =\iota_\gamma^{-1}
        (U_\gamma^*Y_\gamma U_\gamma).
        \label{eq:mpdo_bnt_direct_sum_unitary_S}
    \end{align}
    These are the converse-direction middle conjugations introduced in
    \cite[Appendix~C.4, lines~2053--2080]{Cirac2016MPDO_arXiv}.  The
    existence of the unitaries in
    (\ref{eq:mpdo_bnt_sector_unitary_conjugacy}) is an assumption here.
    They are distinct from the vertical-sector direct-sum maps in
    \cite[Appendix~C.4, lines~1955--1997]{Cirac2016MPDO_arXiv}.  No maps
    between the ambient full matrix algebras are included here.
\end{definition}

\begin{theorem}[Properties of the direct-sum unitary conjugations]
    \label{thm:mpdo_bnt_direct_sum_unitary_conjugations}
    \lean{MPOTensor.BNTAlgebraTensorClause.TwoSiteExactSectorGauge.%
        UnitarySectorConjugacy.directSumUnitaryT_isKrausDirectSumMap}
    \lean{MPOTensor.BNTAlgebraTensorClause.TwoSiteExactSectorGauge.%
        UnitarySectorConjugacy.directSumUnitaryS_isKrausDirectSumMap}
    \lean{MPOTensor.BNTAlgebraTensorClause.TwoSiteExactSectorGauge.%
        UnitarySectorConjugacy.%
        directSumUnitaryT_isTracePreservingBetweenDirectSums}
    \lean{MPOTensor.BNTAlgebraTensorClause.TwoSiteExactSectorGauge.%
        UnitarySectorConjugacy.%
        directSumUnitaryS_isTracePreservingBetweenDirectSums}
    \lean{MPOTensor.BNTAlgebraTensorClause.TwoSiteExactSectorGauge.%
        UnitarySectorConjugacy.directSumUnitaryS_comp_directSumUnitaryT}
    \lean{MPOTensor.BNTAlgebraTensorClause.TwoSiteExactSectorGauge.%
        UnitarySectorConjugacy.directSumUnitaryT_comp_directSumUnitaryS}
    \lean{MPOTensor.BNTAlgebraTensorClause.TwoSiteExactSectorGauge.%
        UnitarySectorConjugacy.directSumUnitaryT_tensor}
    \lean{MPOTensor.BNTAlgebraTensorClause.TwoSiteExactSectorGauge.%
        UnitarySectorConjugacy.directSumUnitaryS_tensor}
    \leanok
    \uses{def:mpdo_bnt_direct_sum_unitary_conjugations}
    The maps $\widetilde T$ and $\widetilde S$ are completely positive and
    preserve the total trace on the direct sums.  They are mutually inverse:
    \begin{align}
        \widetilde S\widetilde T
         & =\id_{\mathcal A_1},
         &
        \widetilde T\widetilde S
         & =\id_{\mathcal A_2^\sigma}.
        \label{eq:mpdo_bnt_direct_sum_unitary_inverses}
    \end{align}
    They also carry the matched tensor letters in both directions:
    \begin{align}
        \widetilde T\bigl((M_\gamma^i)_\gamma\bigr)
         & =\bigl(A_{\sigma(\gamma)}^i\bigr)_\gamma,
        \label{eq:mpdo_bnt_direct_sum_unitary_tensor_letters_fwd}
        \\
        \widetilde S\bigl((A_{\sigma(\gamma)}^i)_\gamma\bigr)
         & =\bigl(M_\gamma^i\bigr)_\gamma.
        \label{eq:mpdo_bnt_direct_sum_unitary_tensor_letters_rev}
    \end{align}
\end{theorem}

\begin{proof}\leanok

    On each summand, conjugation by $U_\gamma$ and the coordinate
    identification $\iota_\gamma$ are completely positive and
    trace-preserving.  Their direct sums therefore have the same properties.
    Unitarity gives
    \begin{align}
        U_\gamma^*U_\gamma
         & =I,
         &
        U_\gamma U_\gamma^*
         & =I,
    \end{align}
    which proves (\ref{eq:mpdo_bnt_direct_sum_unitary_inverses}) sector by
    sector.  Identity~(\ref{eq:mpdo_bnt_direct_sum_unitary_tensor_letters_fwd}) is precisely
    (\ref{eq:mpdo_bnt_sector_unitary_conjugacy}); applying $\widetilde S$ gives
    (\ref{eq:mpdo_bnt_direct_sum_unitary_tensor_letters_rev}).
\end{proof}

\subsection{Ambient sector retractions and normalized embeddings}

\begin{definition}[Ambient sector retractions and normalized embeddings]
    \label{def:mpdo_bnt_ambient_sector_maps}
    \lean{MPOTensor.BNTAlgebraTensorClause.oneSiteAmbientSectorRetraction}
    \lean{MPOTensor.BNTAlgebraTensorClause.oneSiteAmbientSectorRetraction_apply}
    \lean{MPOTensor.BNTAlgebraTensorClause.oneSiteNormalizedAmbientSectorRestoration}
    \lean{MPOTensor.BNTAlgebraTensorClause.oneSiteNormalizedAmbientSectorRestoration_apply}
    \lean{MPOTensor.BNTAlgebraTensorClause.TwoSiteMultiplicitySpectrum.%
        twoSiteAmbientSectorRetraction}
    \lean{MPOTensor.BNTAlgebraTensorClause.TwoSiteMultiplicitySpectrum.%
        twoSiteAmbientSectorRetraction_apply}
    \lean{MPOTensor.BNTAlgebraTensorClause.TwoSiteMultiplicitySpectrum.%
        twoSiteNormalizedAmbientSectorRestoration}
    \lean{MPOTensor.BNTAlgebraTensorClause.TwoSiteMultiplicitySpectrum.%
        twoSiteNormalizedAmbientSectorRestoration_apply}
    \lean{MPOTensor.BNTAlgebraTensorClause.TwoSiteMultiplicitySpectrum.%
        twoSiteNormalizedAmbientSectorRestoration_apply_explicit}
    \lean{MPOTensor.BNTAlgebraTensorClause.TwoSiteMultiplicitySpectrum.%
        relabeledTwoSiteMultiplicityTrace_eq_oneSite}
    \lean{MPOTensor.verticalSectorAmbientRetraction}
    \lean{MPOTensor.verticalSectorAmbientRetraction_apply}
    \lean{MPOTensor.verticalSectorAmbientRetraction_apply_sector}
    \lean{MPOTensor.normalizedVerticalSectorAmbientRestoration}
    \lean{MPOTensor.normalizedVerticalSectorAmbientRestoration_apply}
    \leanok
    \uses{def:mpdo_bnt_two_site_multiplicity_spectrum, def:mpdo_normalized_vertical_sector_maps, def:mpdo_retained_vertical_sector_coordinates, def:mpdo_two_site_blocking, thm:mpdo_bnt_multiplicity_trace_normalization}
    Let
    \begin{align}
        \mathcal A_1
         & =\bigoplus_\gamma \MN{d_\gamma},
         &
        \mathcal A_2^\sigma
         & =\bigoplus_\gamma \MN{\widetilde d_{\sigma(\gamma)}}
        \notag
    \end{align}
    be the one-site sector algebra and the relabelled two-site sector algebra.
    Their retained physical spaces are
    \begin{align}
        \mathcal K_1
         & =\bigoplus_\gamma
        \left(\C^{r_\gamma}\otimes\C^{d_\gamma}\right),
         &
        \mathcal K_2^\sigma
         & =\bigoplus_\gamma
        \left(\C^{\widetilde r_{\sigma(\gamma)}}
        \otimes\C^{\widetilde d_{\sigma(\gamma)}}\right).
        \notag
    \end{align}
    Let $W_1\colon\C^d\to\mathcal K_1$ and
    $W_2\colon\C^{d^2}\to\mathcal K_2^\sigma$ be the corresponding
    coisometries, so that $W_i W_i^\dagger=I$.  Write $\Pi_{1,\gamma}$ for
    the projection onto the $\gamma$-summand of $\mathcal K_1$ and
    $\Pi_{2,\gamma}^\sigma$ for the projection onto the
    $\sigma(\gamma)$-summand of $\mathcal K_2^\sigma$.  Let
    \begin{align}
        \kappa_\ast\colon\MN{d^2}
         & \longrightarrow \MN{d}\otimes\MN{d}
        \notag
    \end{align}
    be the canonical relabelling of a pair of physical indices.  For the
    positive diagonal multiplicity matrices $\mu_\gamma$ and
    $\nu_{\sigma(\gamma)}$, set
    \begin{align}
        m_\gamma
         & =\tr(\mu_\gamma)
        =\tr(\nu_{\sigma(\gamma)}).
        \label{eq:mpdo_bnt_ambient_sector_trace_weights}
    \end{align}
    Define
    \begin{align}
        R_1\colon\MN{d}                  & \longrightarrow\mathcal A_1,
                                         &
        [R_1(X)]_\gamma
                                         & =\tr_{\mathrm{mult}}
        \left(\Pi_{1,\gamma}W_1XW_1^\dagger\Pi_{1,\gamma}\right),
        \label{eq:mpdo_bnt_ambient_sector_R1}
        \\
        R_2\colon\MN{d}\otimes\MN{d}     & \longrightarrow\mathcal A_2^\sigma,
                                         &
        [R_2(Y)]_\gamma
                                         & =\tr_{\mathrm{mult}}
        \left(\Pi_{2,\gamma}^\sigma W_2\kappa_\ast^{-1}(Y)W_2^\dagger
        \Pi_{2,\gamma}^\sigma\right),
        \label{eq:mpdo_bnt_ambient_sector_R2}
        \\
        \widetilde R_1\colon\mathcal A_1 & \longrightarrow\MN{d},
                                         &
        \widetilde R_1((X_\gamma)_\gamma)
                                         & =W_1^\dagger
        \left(\bigoplus_\gamma
        \frac{\mu_\gamma}{m_\gamma}\otimes X_\gamma\right)W_1,
        \label{eq:mpdo_bnt_ambient_sector_R1_tilde}
        \\
        \widetilde R_2\colon\mathcal A_2^\sigma
                                         & \longrightarrow\MN{d}\otimes\MN{d},
                                         &
        \widetilde R_2((Y_\gamma)_\gamma)
                                         & =\kappa_\ast\!\left[
            W_2^\dagger\left(\bigoplus_\gamma
            \frac{\nu_{\sigma(\gamma)}}{m_\gamma}\otimes Y_\gamma\right)W_2
            \right].
        \label{eq:mpdo_bnt_ambient_sector_R2_tilde}
    \end{align}
    In Appendix~C.4 of \cite{Cirac2016MPDO_arXiv}, lines~2067--2068 define
    $R_1$ by~(\ref{eq:mpdo_bnt_ambient_sector_R1}), and lines~2078--2079
    define $R_2$ by~(\ref{eq:mpdo_bnt_ambient_sector_R2}).
    Lines~2081--2083 give the normalized map $\widetilde R_1$ in
    (\ref{eq:mpdo_bnt_ambient_sector_R1_tilde}).  Lines~2058--2071 give
    $\widetilde R_2$, including its common denominator, in
    (\ref{eq:mpdo_bnt_ambient_sector_R2_tilde}).
    These four definitions use only the two vertical decompositions, their
    sector relabelling, and
    (\ref{eq:mpdo_bnt_ambient_sector_trace_weights}).
\end{definition}

\begin{lemma}[Complete positivity of the ambient sector maps]
    \label{thm:mpdo_bnt_ambient_sector_complete_positivity}
    \lean{MPOTensor.BNTAlgebraTensorClause.oneSiteAmbientSectorRetractionExtension_isKrausCP}
    \lean{MPOTensor.BNTAlgebraTensorClause.oneSiteAmbientSectorRetractionExtension}
    \lean{MPOTensor.BNTAlgebraTensorClause.%
        oneSiteNormalizedAmbientSectorRestorationExtension_isKrausCP}
    \lean{MPOTensor.BNTAlgebraTensorClause.oneSiteNormalizedAmbientSectorRestorationExtension}
    \lean{MPOTensor.BNTAlgebraTensorClause.TwoSiteMultiplicitySpectrum.%
        twoSiteAmbientSectorRetractionExtension_isKrausCP}
    \lean{MPOTensor.BNTAlgebraTensorClause.TwoSiteMultiplicitySpectrum.%
        twoSiteAmbientSectorRetractionExtension}
    \lean{MPOTensor.BNTAlgebraTensorClause.TwoSiteMultiplicitySpectrum.%
        twoSiteNormalizedAmbientSectorRestorationExtension_isKrausCP}
    \lean{MPOTensor.BNTAlgebraTensorClause.TwoSiteMultiplicitySpectrum.%
        twoSiteNormalizedAmbientSectorRestorationExtension}
    \lean{MPOTensor.verticalSectorAmbientRetractionExtension}
    \lean{MPOTensor.verticalSectorAmbientRetractionExtension_factorization}
    \lean{MPOTensor.verticalSectorAmbientRetractionExtension_isKrausCP}
    \lean{MPOTensor.normalizedVerticalSectorAmbientRestorationExtension}
    \lean{MPOTensor.normalizedVerticalSectorAmbientRestorationExtension_factorization}
    \lean{MPOTensor.normalizedVerticalSectorAmbientRestorationExtension_isKrausCP}
    \leanok
    \uses{def:mpdo_bnt_ambient_sector_maps}
    When each sector family is represented by its block-diagonal matrix and
    each map on a sector algebra first extracts the diagonal sector blocks
    from an arbitrary matrix, the resulting canonical full-matrix
    representatives of $R_1$, $R_2$, $\widetilde R_1$, and $\widetilde R_2$
    are completely positive.
\end{lemma}

\begin{proof}\leanok
    \uses{thm:mpdo_vertical_sector_partial_trace_kraus}
    Each retraction $R_i$ is the composition of compression by $W_i$, a unitary
    relabelling of coordinates, sectorwise partial trace, and block-diagonal
    inclusion.  Each normalized embedding is the composition of diagonal-block
    extraction, preparation of the positive matrix
    $\mu_\gamma/m_\gamma$ or
    $\nu_{\sigma(\gamma)}/m_\gamma$, coordinate relabelling, and expansion by
    $W_i^\dagger$.  Every constituent has a Kraus representation, and
    compositions of such maps are completely positive.
\end{proof}

% Local fix (zero-sector complement): the coisometries need not have full
% ambient support, so the raw maps are trace-nonincreasing rather than
% trace-preserving. Documented in
% docs/paper-gaps/cpgsv17_vertical_isometry_zero_sector.tex.
\begin{lemma}[Trace of the ambient sector maps]
    \label{thm:mpdo_bnt_ambient_sector_trace}
    \lean{MPOTensor.BNTAlgebraTensorClause.verticalSectorTrace_oneSiteAmbientSectorRetraction}
    \lean{MPOTensor.BNTAlgebraTensorClause.TwoSiteMultiplicitySpectrum.%
        verticalSectorTrace_twoSiteAmbientSectorRetraction}
    \lean{MPOTensor.BNTAlgebraTensorClause.verticalSectorTrace_oneSiteAmbientSectorRetraction_le}
    \lean{MPOTensor.BNTAlgebraTensorClause.TwoSiteMultiplicitySpectrum.%
        verticalSectorTrace_twoSiteAmbientSectorRetraction_le}
    \lean{MPOTensor.BNTAlgebraTensorClause.trace_oneSiteNormalizedAmbientSectorRestoration}
    \lean{MPOTensor.BNTAlgebraTensorClause.TwoSiteMultiplicitySpectrum.%
        trace_twoSiteNormalizedAmbientSectorRestoration}
    \lean{MPOTensor.BNTAlgebraTensorClause.%
        oneSiteNormalizedAmbientSectorRestorationExtension_isKrausCPTP}
    \lean{MPOTensor.BNTAlgebraTensorClause.TwoSiteMultiplicitySpectrum.%
        twoSiteNormalizedAmbientSectorRestorationExtension_isKrausCPTP}
    \lean{MPOTensor.verticalSectorTrace_verticalSectorAmbientRetraction_le}
    \lean{MPOTensor.trace_normalizedVerticalSectorAmbientRestoration}
    \lean{MPOTensor.normalizedVerticalSectorAmbientRestorationExtension_isKrausCPTP}
    \leanok
    \uses{def:mpdo_bnt_ambient_sector_maps, def:mpdo_positive_vertical_sector_family, thm:mpdo_bnt_ambient_sector_complete_positivity}
    For arbitrary ambient matrices $X\in\MN{d}$ and
    $Y\in\MN{d}\otimes\MN{d}$,
    \begin{align}
        \Trsec(R_1(X))
         & =\tr(W_1XW_1^\dagger),
        \label{eq:mpdo_bnt_ambient_sector_R1_trace}
        \\
        \Trsec(R_2(Y))
         & =\tr\!\left(W_2\kappa_\ast^{-1}(Y)W_2^\dagger\right).
        \label{eq:mpdo_bnt_ambient_sector_R2_trace}
    \end{align}
    If $X$ and $Y$ are positive semidefinite, then
    \begin{align}
        \Trsec(R_1(X))
         & \leq\tr(X),
         &
        \Trsec(R_2(Y))
         & \leq\tr(Y).
        \label{eq:mpdo_bnt_ambient_sector_trace_loss}
    \end{align}
    In the other direction, for $Z\in\mathcal A_1$ and
    $Z'\in\mathcal A_2^\sigma$,
    \begin{align}
        \tr(\widetilde R_1(Z))
         & =\Trsec(Z),
         &
        \tr(\widetilde R_2(Z'))
         & =\Trsec(Z').
        \label{eq:mpdo_bnt_ambient_sector_embedding_trace}
    \end{align}
    Hence the canonical full-matrix representatives of $\widetilde R_1$ and
    $\widetilde R_2$ are trace-preserving completely positive maps.

    The inequalities in (\ref{eq:mpdo_bnt_ambient_sector_trace_loss}) may be
    strict: $W_i^\dagger W_i$ projects onto the retained nonzero sectors, so
    positive elements supported on the discarded orthogonal complement lose
    trace under the compression.
\end{lemma}

\begin{proof}\leanok
    \uses{lem:mpdo_normalized_vertical_sector_positive_trace, lem:mpdo_coisometric_compression_trace}
    Sectorwise partial trace and block-diagonal inclusion preserve the total
    trace, which gives
    (\ref{eq:mpdo_bnt_ambient_sector_R1_trace}) and
    (\ref{eq:mpdo_bnt_ambient_sector_R2_trace}).  For
    $P_i=W_i^\dagger W_i$, cyclicity of trace gives
    \begin{align}
        \tr(W_i ZW_i^\dagger)
         & =\tr(P_i Z)\leq\tr(Z)
        \notag
    \end{align}
    when $Z\succeq0$.  The relabelling $\kappa_\ast$ preserves positivity and
    trace, proving (\ref{eq:mpdo_bnt_ambient_sector_trace_loss}).  Finally, each
    normalized multiplicity matrix has trace one, and
    $W_i W_i^\dagger=I$ gives
    $\tr(W_i^\dagger BW_i)=\tr(B)$.  This proves
    (\ref{eq:mpdo_bnt_ambient_sector_embedding_trace}); complete positivity
    follows from Lemma~\ref{thm:mpdo_bnt_ambient_sector_complete_positivity}.
\end{proof}

\begin{lemma}[Retraction identities for the ambient sector maps]
    \label{thm:mpdo_bnt_ambient_sector_retractions}
    \lean{MPOTensor.BNTAlgebraTensorClause.%
        oneSiteAmbientSectorRetraction_comp_normalizedRestoration}
    \lean{MPOTensor.BNTAlgebraTensorClause.TwoSiteMultiplicitySpectrum.%
        twoSiteAmbientSectorRetraction_comp_normalizedRestoration}
    \lean{MPOTensor.verticalSectorAmbientRetraction_comp_normalizedRestoration}
    \leanok
    \uses{def:mpdo_bnt_ambient_sector_maps}
    The retractions $R_1$ and $R_2$ are left inverses of their normalized embeddings:
    \begin{align}
        R_1\widetilde R_1
         & =\id_{\mathcal A_1},
         &
        R_2\widetilde R_2
         & =\id_{\mathcal A_2^\sigma}.
        \label{eq:mpdo_bnt_ambient_sector_retractions}
    \end{align}
\end{lemma}

\begin{proof}\leanok
    \uses{lem:mpdo_retained_vertical_sector_retraction}
    Compressing an expansion by $W_i^\dagger$ gives the original retained
    block-diagonal matrix because $W_i W_i^\dagger=I$.  On every sector,
    \begin{align}
        \tr_{\mathrm{mult}}
        \left(\frac{\mu_\gamma}{m_\gamma}\otimes X_\gamma\right)
         & =X_\gamma,
        \notag        \\
        \tr_{\mathrm{mult}}
        \left(\frac{\nu_{\sigma(\gamma)}}{m_\gamma}\otimes Y_\gamma\right)
         & =Y_\gamma,
        \label{eq:mpdo_bnt_ambient_retract_nu}
    \end{align}
    where (\ref{eq:mpdo_bnt_ambient_retract_nu}) follows from
    $\tr(\nu_{\sigma(\gamma)})=m_\gamma$.  The two coordinate relabellings in
    the two-site composite cancel.  These identities prove
    (\ref{eq:mpdo_bnt_ambient_sector_retractions}).
\end{proof}

\begin{lemma}[Action on the weighted tensor letters]
    \label{thm:mpdo_bnt_ambient_sector_tensor_letters}
    \lean{MPOTensor.BNTAlgebraTensorClause.oneSiteAmbientSectorRetraction_verticalTensor}
    \lean{MPOTensor.BNTAlgebraTensorClause.%
        oneSiteNormalizedAmbientSectorRestoration_trace_smul_verticalTensor}
    \lean{MPOTensor.BNTAlgebraTensorClause.TwoSiteMultiplicitySpectrum.%
        twoSiteAmbientSectorRetraction_verticalTensor}
    \lean{MPOTensor.BNTAlgebraTensorClause.TwoSiteMultiplicitySpectrum.%
        twoSiteNormalizedAmbientSectorRestoration_trace_smul_verticalTensor}
    \lean{MPOTensor.verticalSectorAmbientRetraction_weightedBlockDiagonal}
    \lean{MPOTensor.normalizedVerticalSectorAmbientRestoration_trace_smul}
    \leanok
    \uses{def:mpdo_bnt_ambient_sector_maps, lem:mpdo_normalized_retained_vertical_sector_embedding, lem:mpdo_retained_vertical_sector_partial_trace, thm:mpdo_bnt_multiplicity_trace_normalization}
    For every vertical letter $v$, the one-site maps satisfy
    \begin{align}
        R_1(\widetilde M_v)
         & =\bigl(m_\gamma M_\gamma^v\bigr)_\gamma,
        \label{eq:mpdo_bnt_ambient_sector_one_site_fwd}
        \\
        \widetilde R_1\bigl((m_\gamma M_\gamma^v)_\gamma\bigr)
         & =\widetilde M_v.
        \label{eq:mpdo_bnt_ambient_sector_one_site_rev}
    \end{align}
    The relabelled two-site maps satisfy
    \begin{align}
        R_2\!\left(\kappa_\ast(\widetilde{M^{[2]}}_v)\right)
         & =\bigl(m_\gamma A_{\sigma(\gamma)}^v\bigr)_\gamma,
        \label{eq:mpdo_bnt_ambient_sector_two_site_fwd}
        \\
        \widetilde R_2\bigl((m_\gamma A_{\sigma(\gamma)}^v)_\gamma\bigr)
         & =\kappa_\ast(\widetilde{M^{[2]}}_v).
        \label{eq:mpdo_bnt_ambient_sector_two_site_rev}
    \end{align}
    Identities~(\ref{eq:mpdo_bnt_ambient_sector_one_site_fwd}) and~(\ref{eq:mpdo_bnt_ambient_sector_one_site_rev}) combine the one-site weighted
    sector form of
    \cite[Appendix~C.4, lines~2048--2051]{Cirac2016MPDO_arXiv} with,
    respectively,
    \cite[Appendix~C.4, lines~2067--2068]{Cirac2016MPDO_arXiv} and
    \cite[Appendix~C.4, lines~2081--2083]{Cirac2016MPDO_arXiv}.
    Identity~(\ref{eq:mpdo_bnt_ambient_sector_two_site_fwd}) combines the matched weighted sector form of
    \cite[Appendix~C.4, lines~2048--2064]{Cirac2016MPDO_arXiv} with the map
    $R_2$ of \cite[Appendix~C.4, lines~2078--2079]{Cirac2016MPDO_arXiv}.
    Identity~(\ref{eq:mpdo_bnt_ambient_sector_two_site_rev}) combines the same sector form with the normalized map
    $\widetilde R_2$ of
    \cite[Appendix~C.4, lines~2070--2071]{Cirac2016MPDO_arXiv}.
\end{lemma}

\begin{proof}\leanok
    \uses{lem:mpdo_normalized_retained_vertical_sector_embedding, lem:mpdo_retained_vertical_sector_partial_trace, thm:mpdo_bnt_multiplicity_trace_normalization}
    The vertical reconstruction identities give
    \begin{align}
        \widetilde M_v
         & =W_1^\dagger\left(\bigoplus_\gamma
        \mu_\gamma\otimes M_\gamma^v\right)W_1,
        \notag                                \\
        \widetilde{M^{[2]}}_v
         & =W_2^\dagger\left(\bigoplus_\gamma
        \nu_{\sigma(\gamma)}\otimes A_{\sigma(\gamma)}^v\right)W_2.
        \notag
    \end{align}
    Compression removes the two outer coisometries, and partial trace over the
    multiplicity coordinate multiplies the simple-sector letter by
    $\tr(\mu_\gamma)=m_\gamma$ or
    $\tr(\nu_{\sigma(\gamma)})=m_\gamma$.  This proves
    (\ref{eq:mpdo_bnt_ambient_sector_one_site_fwd}) and
    (\ref{eq:mpdo_bnt_ambient_sector_two_site_fwd}).  Preparing the
    normalized multiplicity matrix cancels the factor $m_\gamma$, yielding
    (\ref{eq:mpdo_bnt_ambient_sector_one_site_rev}) and
    (\ref{eq:mpdo_bnt_ambient_sector_two_site_rev}).
\end{proof}

\subsection{Trace-preserving completions and physical channels}

\begin{theorem}[Trace-preserving completion outside a projection]
    \label{thm:matrix_support_completion_cptp}
    \lean{Matrix.supportCompletion_isKrausCPTP}
    \lean{Matrix.tracePrepareMap_trace}
    \lean{Matrix.tracePrepareMap_isKrausCP}
    \lean{Matrix.tracePrepareMap_isKrausCPTP}
    \lean{Matrix.supportComplementMap}
    \lean{Matrix.supportComplementMap_apply}
    \lean{Matrix.supportComplementMap_isKrausCP}
    \lean{Matrix.supportComplementMap_trace}
    \lean{Matrix.supportCompletion}
    \lean{Matrix.supportCompletion_apply}
    \lean{Matrix.supportCompletion_isKrausCP}
    \lean{Matrix.supportCompletion_trace}
    \lean{Matrix.supportCompletion_apply_of_supported}
    \leanok

    Let $P\in\MN{m}$ be an orthogonal projection, let
    $\rho\in\MN{n}$ be positive semidefinite with $\tr(\rho)=1$, and let
    $\mathcal E\colon\MN{m}\to\MN{n}$ be completely positive.  Define the
    trace-and-prepare map and its complementary restriction by
    \begin{align}
        \mathcal D_\rho(X)
         & =\tr(X)\rho,
         &
        \mathcal C_{P,\rho}(X)
         & =\mathcal D_\rho\bigl((I-P)X(I-P)^\dagger\bigr).
        \label{eq:matrix_support_complement_map}
    \end{align}
    Both maps are completely positive, and $\mathcal D_\rho$ is
    trace-preserving.  If
    \begin{align}
        \tr(\mathcal E(X))=\tr(PX)
        \label{eq:matrix_support_map_trace}
    \end{align}
    for every $X\in\MN{m}$, then the completion
    \begin{align}
        \widehat{\mathcal E}(X)
         & =\mathcal E(X)+\mathcal C_{P,\rho}(X)
        \label{eq:matrix_support_completion}
    \end{align}
    is trace-preserving and completely positive.  It agrees with
    $\mathcal E$ on every matrix supported by $P$:
    \begin{align}
        PXP=X
        \quad\Longrightarrow\quad
        \widehat{\mathcal E}(X)=\mathcal E(X).
        \label{eq:matrix_support_completion_agreement}
    \end{align}
\end{theorem}

\begin{proof}\leanok

    The map $\mathcal D_\rho$ factors as state preparation, interchange of
    the two tensor factors, and partial trace.  The map
    $X\mapsto(I-P)X(I-P)^\dagger$ has a single Kraus operator.  Thus the two
    maps in~(\ref{eq:matrix_support_complement_map}) are completely positive.
    Since $P=P^\dagger=P^2$,
    \begin{align}
        \tr\bigl(\mathcal C_{P,\rho}(X)\bigr)
         & =\tr((I-P)X).
        \notag
    \end{align}
    Adding this identity to~(\ref{eq:matrix_support_map_trace}) gives
    $\tr(\widehat{\mathcal E}(X))=\tr(X)$.  Finally, $PXP=X$ implies
    $(I-P)X(I-P)=0$, which proves
    (\ref{eq:matrix_support_completion_agreement}).
\end{proof}

\begin{theorem}[Physical channels from unitary sector conjugacies]
    \label{thm:mpdo_bnt_conditional_physical_channels}
    \lean{MPOTensor.BNTAlgebraTensorClause.TwoSiteExactSectorGauge.%
        UnitarySectorConjugacy.isRFPViaTS}
    \lean{MPOTensor.BNTAlgebraTensorClause.oneSiteRetainedProjection}
    \lean{MPOTensor.BNTAlgebraTensorClause.oneSiteRetainedProjection_isHermitian}
    \lean{MPOTensor.BNTAlgebraTensorClause.oneSiteRetainedProjection_mul_self}
    \lean{MPOTensor.BNTAlgebraTensorClause.TwoSiteMultiplicitySpectrum.%
        twoSiteRetainedProjection}
    \lean{MPOTensor.BNTAlgebraTensorClause.TwoSiteMultiplicitySpectrum.%
        twoSiteRetainedProjection_isHermitian}
    \lean{MPOTensor.BNTAlgebraTensorClause.TwoSiteMultiplicitySpectrum.%
        twoSiteRetainedProjection_mul_self}
    \lean{MPOTensor.BNTAlgebraTensorClause.TwoSiteExactSectorGauge.%
        UnitarySectorConjugacy.rawPhysicalT}
    \lean{MPOTensor.BNTAlgebraTensorClause.TwoSiteExactSectorGauge.%
        UnitarySectorConjugacy.rawPhysicalS}
    \lean{MPOTensor.BNTAlgebraTensorClause.TwoSiteExactSectorGauge.%
        UnitarySectorConjugacy.rawPhysicalT_isKrausCP}
    \lean{MPOTensor.BNTAlgebraTensorClause.TwoSiteExactSectorGauge.%
        UnitarySectorConjugacy.rawPhysicalS_isKrausCP}
    \lean{MPOTensor.BNTAlgebraTensorClause.TwoSiteExactSectorGauge.%
        UnitarySectorConjugacy.trace_rawPhysicalT}
    \lean{MPOTensor.BNTAlgebraTensorClause.TwoSiteExactSectorGauge.%
        UnitarySectorConjugacy.trace_rawPhysicalS}
    \lean{MPOTensor.BNTAlgebraTensorClause.TwoSiteExactSectorGauge.%
        UnitarySectorConjugacy.rawPhysicalT_physClose1}
    \lean{MPOTensor.BNTAlgebraTensorClause.TwoSiteExactSectorGauge.%
        UnitarySectorConjugacy.rawPhysicalS_physClose2}
    \lean{MPOTensor.BNTAlgebraTensorClause.TwoSiteExactSectorGauge.%
        oneSiteRetainedProjection_physClose1}
    \lean{MPOTensor.BNTAlgebraTensorClause.TwoSiteExactSectorGauge.%
        twoSiteRetainedProjection_physClose2}
    \lean{MPOTensor.BNTAlgebraTensorClause.TwoSiteExactSectorGauge.%
        UnitarySectorConjugacy.physicalT}
    \lean{MPOTensor.BNTAlgebraTensorClause.TwoSiteExactSectorGauge.%
        UnitarySectorConjugacy.physicalS}
    \lean{MPOTensor.BNTAlgebraTensorClause.TwoSiteExactSectorGauge.%
        UnitarySectorConjugacy.physicalT_isKrausCPTP}
    \lean{MPOTensor.BNTAlgebraTensorClause.TwoSiteExactSectorGauge.%
        UnitarySectorConjugacy.physicalS_isKrausCPTP}
    \lean{MPOTensor.BNTAlgebraTensorClause.TwoSiteExactSectorGauge.%
        UnitarySectorConjugacy.physicalT_physClose1}
    \lean{MPOTensor.BNTAlgebraTensorClause.TwoSiteExactSectorGauge.%
        UnitarySectorConjugacy.physicalS_physClose2}
    \leanok
    \uses{def:mpdo_bnt_ambient_sector_maps, def:mpdo_bnt_direct_sum_unitary_conjugations, def:mpdo_faithful_uniform_density, def:is_rfp_via_ts}
    Suppose that the one-site and two-site sector decompositions are paired by
    a relabelling $\sigma$ and unitaries $U_\gamma$ satisfying
    \begin{align}
        A_{\sigma(\gamma)}^i
         & =U_\gamma M_\gamma^i U_\gamma^\dagger
        \label{eq:mpdo_bnt_physical_channel_unitary_conjugacy}
    \end{align}
    for every sector $\gamma$ and every vertical letter $i$.  With the ambient
    sector maps and unitary conjugations defined above, set
    \begin{align}
        \mathcal T_0
         & =\widetilde R_2\circ\widetilde T\circ R_1,
         &
        \mathcal S_0
         & =\widetilde R_1\circ\widetilde S\circ R_2.
        \label{eq:mpdo_bnt_active_physical_maps}
    \end{align}
    Let
    \begin{align}
        P_1 & =W_1^\dagger W_1,
            &
        P_2 & =\kappa_\ast(W_2^\dagger W_2),
            &
        Q_i & =I-P_i.
        \label{eq:mpdo_bnt_physical_active_projections}
    \end{align}
    The matrices $P_1$ and $P_2$ are orthogonal projections.  If $d>0$, put
    $\tau_1=d^{-1}I_d$ and $\tau_2=d^{-2}I_{d^2}$, and define
    \begin{align}
        \mathcal T(X)
         & =\mathcal T_0(X)+\tr(Q_1XQ_1)\tau_2,
        \notag                                  \\
        \mathcal S(Y)
         & =\mathcal S_0(Y)+\tr(Q_2YQ_2)\tau_1.
        \label{eq:mpdo_bnt_completed_physical_maps}
    \end{align}
    If $d=0$, take both maps to be zero; their domain and codomain matrix
    algebras then consist only of zero.
    Then $\mathcal T$ and $\mathcal S$ are trace-preserving completely positive
    maps on the full physical matrix algebras.  For every virtual matrix $Z$,
    they satisfy
    \begin{align}
        \mathcal T[M_1(Z)] & =M_2(Z),
                           &
        \mathcal S[M_2(Z)] & =M_1(Z).
        \label{eq:mpdo_bnt_completed_physical_map_actions}
    \end{align}
    Hence $M$ is a renormalization fixed point in the sense of
    Definition~\ref{def:is_rfp_via_ts}.

    The compositions in
    (\ref{eq:mpdo_bnt_active_physical_maps}) are the formulas of
    \cite[Appendix~C.4, lines~2065--2085]{Cirac2016MPDO_arXiv}.  Those formulas
    determine the action after compression to the retained physical sectors.
    With rectangular coisometries, the uncompleted composites annihilate the
    discarded zero-sector complements and fail to preserve trace there;
    Formula~(\ref{eq:mpdo_bnt_completed_physical_maps}) provides a
    trace-preserving completion on those complements when the physical space is
    nonzero, yielding the same action on the physical operators generated by $M$
    for any choice of reference states $\tau_1,\tau_2$.
\end{theorem}

\begin{proof}\leanok
    \uses{thm:mpdo_bnt_ambient_sector_complete_positivity, thm:mpdo_bnt_ambient_sector_trace, thm:mpdo_bnt_ambient_sector_tensor_letters, thm:mpdo_bnt_direct_sum_unitary_conjugations, thm:matrix_support_completion_cptp, thm:mpdo_faithful_uniform_density, lem:mpdo_coisometric_compression_trace, lem:mpdo_vertical_bond_contraction_phys_close, lem:mpdo_vertical_bond_contraction_reconstruction, thm:mpdo_one_site_closure_two_site_blocking}
    Complete positivity of the four outer maps and of the two direct-sum
    conjugations gives complete positivity of $\mathcal T_0$ and
    $\mathcal S_0$.  Their traces are
    \begin{align}
        \tr(\mathcal T_0(X)) & =\tr(P_1X),
                             &
        \tr(\mathcal S_0(Y)) & =\tr(P_2Y).
        \label{eq:mpdo_bnt_active_physical_map_traces}
    \end{align}
    Suppose first that $d>0$.  Since $Q_i$ is a projection and $\tau_i$ has
    trace one, the corresponding complementary term is completely positive
    and has trace $\tr(Q_iXQ_i)=\tr(Q_iX)$.  Adding
    (\ref{eq:mpdo_bnt_active_physical_map_traces}) and using $P_i+Q_i=I$
    proves trace preservation of the completed maps.  If $d=0$, every ambient
    matrix is zero, so the zero maps are completely positive and trace
    preserving.

    The reconstruction identities for the two vertical decompositions imply
    \begin{align}
        P_1M_1(Z)P_1 & =M_1(Z),
                     &
        P_2M_2(Z)P_2 & =M_2(Z).
        \label{eq:mpdo_bnt_physical_closure_support}
    \end{align}
    Thus both complementary terms vanish on the corresponding physical
    operators.  The retraction formulas extract the families
    $(m_\gamma M_\gamma^i)_\gamma$ and
    $(m_\gamma A_{\sigma(\gamma)}^i)_\gamma$; unitary conjugation carries
    these families into one another, and the normalized restorations recover
    the one-site and two-site letters.  Linearity in the virtual matrix gives
    (\ref{eq:mpdo_bnt_completed_physical_map_actions}).
\end{proof}

\begin{theorem}[BNT algebra clause implies a renormalization fixed point]
    \label{thm:mpdo_bnt_algebra_tensor_to_rfp}
    \lean{MPOTensor.BNTAlgebraTensorClause.isRFPViaTS}
    \lean{MPOTensor.HasBNTAlgebraTensorClause.isRFPViaTS}
    \leanok
    \uses{def:mpdo, def:cpsv_canonical_form, def:mpdo_bnt_algebra_tensor_clause, def:is_rfp_via_ts}
    Let $M$ be an MPDO whose doubled-index tensor is in CPSV
    canonical form.  As in implication (ii)$\Rightarrow$(i) of
    \cite[Theorem~4.14, lines~972--993]{Cirac2016MPDO_arXiv}, if $M$ has a
    tensor-attached basis-of-normal-tensors algebra clause, then there are
    trace-preserving completely positive maps
    \begin{align}
        \mathcal T\colon\MN d    & \longrightarrow\MN{d^2},
                                 &
        \mathcal S\colon\MN{d^2} & \longrightarrow\MN d
        \notag
    \end{align}
    such that, for every virtual matrix $X$,
    \begin{align}
        \mathcal T[M_1(X)] & =M_2(X),
                           &
        \mathcal S[M_2(X)] & =M_1(X).
        \label{eq:mpdo_bnt_algebra_tensor_rfp_maps}
    \end{align}
    Thus $M$ is a renormalization fixed point in the sense of
    Definition~\ref{def:is_rfp_via_ts}.

    Here the common marked comparison needed to normalize the exact sector
    gauges comes from
    Theorem~\ref{thm:mpdo_bnt_mixed_prefix_unitary_sector_conjugacy}
    \cite{gap:cpsv16_two_site_sector_unitary_gauge_gap}, with the physical maps
    completed on the discarded complements as in
    \cite{gap:cpgsv17_vertical_isometry_zero_sector}.
\end{theorem}

\begin{proof}\leanok
    \uses{thm:mpdo_bnt_algebra_tensor_clause_spectrum, thm:mpdo_bnt_mixed_prefix_unitary_sector_conjugacy, thm:mpdo_bnt_conditional_physical_channels}
    The tensor-attached algebra clause determines the two-site vertical
    decomposition, the sector relabelling, and exact invertible gauges, and
    Theorem~\ref{thm:mpdo_bnt_mixed_prefix_unitary_sector_conjugacy}
    normalizes these gauges simultaneously to unitaries, so
    Theorem~\ref{thm:mpdo_bnt_conditional_physical_channels} yields the two
    trace-preserving completely positive maps and
    (\ref{eq:mpdo_bnt_algebra_tensor_rfp_maps}).
\end{proof}



